% ============================================================
%  Biodiversity and Ecosystems — Experiment Booklet (one section per experiment)
%  Year 7 (Year 8 extension)  |  Victorian Curriculum F–10 Version 2.0
%  Compile with: xelatex three times, then python3 fill.py Ecosystems_Experiment_Booklet_Yr7
% ============================================================
\documentclass[11pt,a4paper]{article}
\newcommand{\unitname}{Biodiversity and Ecosystems}\newcommand{\unitcaps}{BIODIVERSITY AND ECOSYSTEMS}

% ── Fonts (compile with LuaLaTeX) ────────────────────────────────────────────
\usepackage{fontspec}
% Bundled fonts (in ./fonts/) so the document compiles anywhere with XeLaTeX/LuaLaTeX
\setmainfont{Lato}[                   % body — humanist sans (stands in for deck's Inter)
  Path=./fonts/, Extension=.ttf, Ligatures=TeX,
  UprightFont=Lato-Regular, ItalicFont=Lato-Italic,
  BoldFont=Lato-Bold, BoldItalicFont=Lato-BoldItalic ]
\newfontfamily\headingfont{Poppins}[  % display/headings — matches the slide decks
  Path=./fonts/, Extension=.ttf, Ligatures=TeX,
  UprightFont=Poppins-Regular, ItalicFont=Poppins-Italic,
  BoldFont=Poppins-Bold, BoldItalicFont=Poppins-BoldItalic ]
\setlength{\headheight}{22pt}

\usepackage[protrusion=true,expansion=false]{microtype}
\usepackage[a4paper, top=2cm, bottom=2cm, left=2cm, right=2cm]{geometry}
\usepackage{xcolor}
\usepackage{tikz}
\usetikzlibrary{arrows.meta, shapes.geometric, calc, positioning, decorations.pathreplacing}
\usepackage{tcolorbox}
\tcbuselibrary{skins, breakable, raster}
\usepackage{fancyhdr}
\usepackage{titlesec}
\usepackage{enumitem}
\usepackage{booktabs}
\usepackage{tabularx}
\usepackage{array}
\usepackage{colortbl}
\usepackage{multicol}
\usepackage{amsmath}
\usepackage{needspace}
% TikZ diagram macros are embedded below (just before \begin{document}) for easy editing.
\usepackage{longtable}
\usepackage{lastpage}
\usepackage{calc}
\usepackage{mdframed}
\usepackage{setspace}
\usepackage{ragged2e}
% PDF bookmarks: one per lesson (made in \sectiondiv) plus the cover and
% contents, and links from the contents table to each lesson.
% implicit=false: hyperref adds no anchors of its own (an anchor on each
% list item nudged lists down 1.5 mm, which would undo the page fill).
% hidelinks keeps the links invisible, so the printed pages are unchanged.
\usepackage[implicit=false, hidelinks,
  pdftitle={Biodiversity and Ecosystems — Experiment Booklet},
  pdfsubject={Years 7–8 Science, Victorian Curriculum F–10 v2.0}]{hyperref}
\usepackage{bookmark} % writes the outline itself, which implicit=false switches off in hyperref
% hyperref 7.01 with implicit=false leaves \Hy@PageAnchorSlide undefined, which
% throws at \end{document} and stops fill.py (it builds with -halt-on-error).
\makeatletter\providecommand\Hy@PageAnchorSlide{}\makeatother

% ── Colour Palette ───────────────────────────────────────────────────────────
% Remapped to the slide-deck palette (names kept so body + diagrams cascade)
\definecolor{NavyBlue}   {HTML}{16305C}
\definecolor{MidBlue}    {HTML}{2E6FE0}
\definecolor{LightBlue}  {HTML}{E8F1FF}
\definecolor{TealGreen}  {HTML}{0F6E63}
\definecolor{LightGreen} {HTML}{E0F7F3}
\definecolor{GoldDark}   {HTML}{C77F12}
\definecolor{GoldLight}  {HTML}{FFF3DC}
\definecolor{PurpleDark} {HTML}{5B3FC0}
\definecolor{PurpleLight}{HTML}{EFE9FF}
\definecolor{OrangeDark} {HTML}{C9483D}
\definecolor{OrangeLight}{HTML}{FFE9E6}
\definecolor{MidGrey}    {HTML}{AEBED4}
\definecolor{LightGrey}  {HTML}{F2F5FA}
\definecolor{DarkGrey}   {HTML}{3A4759}
\definecolor{TableHead}  {HTML}{16305C}
\definecolor{RowAlt}     {HTML}{EEF4FF}
% Bright accent tints (deck) for bars, rules and chips
\definecolor{Blue}   {HTML}{2E6FE0}
\definecolor{Teal}   {HTML}{18A999}
\definecolor{Amber}  {HTML}{F5A623}
\definecolor{Coral}  {HTML}{FF6B5E}
\definecolor{Violet} {HTML}{7C5CE0}
\definecolor{Sky}    {HTML}{4DA3FF}
\definecolor{Ink}    {HTML}{16223A}
\definecolor{Line}   {HTML}{DDE5F0}
\definecolor{Paper}  {HTML}{FBFCFE}

% ── Page Style ───────────────────────────────────────────────────────────────
\pagestyle{fancy}
\fancyhf{}
\renewcommand{\headrulewidth}{0.4pt}
\renewcommand{\headrule}{\hbox to\headwidth{\color{NavyBlue}\leaders\hrule height \headrulewidth\hfill}}
\fancyhead[L]{\headingfont\small\color{NavyBlue}\unitname{} \textbar{} Experiment Booklet}
\fancyhead[R]{\headingfont\small\color{DarkGrey}Year 7 \textbar{} Victorian Curriculum F--10 v2.0}
\fancyfoot[C]{\headingfont\small\color{DarkGrey}\thepage\ \textbar{} \pageref{LastPage}}
\fancypagestyle{plain}{%
  \fancyhf{}
  \fancyfoot[C]{\headingfont\small\color{DarkGrey}\thepage\ \textbar{} \pageref{LastPage}}
  \renewcommand{\headrulewidth}{0pt}
}

% ── Section Formatting ───────────────────────────────────────────────────────
\titleformat{\section}
  {\headingfont\color{NavyBlue}\Large\bfseries}
  {}{0em}{}
  [{\color{Teal}\rule[0.6ex]{2.2em}{3pt}}{\color{Line}\rule[0.6ex]{\dimexpr\linewidth-2.2em\relax}{1pt}}]
\titlespacing{\section}{0pt}{18pt}{8pt}

\titleformat{\subsection}
  {\headingfont\color{MidBlue}\large\bfseries}
  {\thesubsection}{0.6em}{}
\titlespacing{\subsection}{0pt}{12pt}{4pt}

\titleformat{\subsubsection}
  {\headingfont\color{NavyBlue}\normalsize\bfseries}
  {}{0em}{}
\titlespacing{\subsubsection}{0pt}{8pt}{2pt}

% Deck-style eyebrow / kicker (optional use)
\newcommand{\kicker}[2][Teal]{%
  {\headingfont\bfseries\footnotesize\color{#1}%
   \textcolor{#1}{\rule[0.18ex]{1.6em}{2.2pt}}\hspace{0.5em}\MakeUppercase{#2}\par}%
  \vspace{2pt}}

% ── Custom tcolorbox Environments ────────────────────────────────────────────

% Section divider banner
\newtcolorbox{sectionbanner}[2]{
  enhanced, breakable=false,
  colback=NavyBlue, colframe=NavyBlue,
  fontupper=\color{white}\bfseries\large,
  left=4mm, right=4mm, top=3mm, bottom=3mm,
  title={\color{white}\Huge\bfseries #1},
  attach boxed title to top left={yshift=-2mm, xshift=4mm},
  boxed title style={colback=NavyBlue, colframe=NavyBlue},
  after upper={\par\smallskip\normalfont\normalsize\color{white!80!NavyBlue}\textit{#2}}
}

% Key info box (blue)
\newtcolorbox{infobox}[1]{
  breakable, arc=2.5mm, boxrule=0.5pt,
  colback=LightBlue, colframe=Blue!40,
  fonttitle=\headingfont\bfseries, coltitle=MidBlue,
  title={#1},
  left=4mm, right=3.5mm, top=2.5mm, bottom=2.5mm
}

% Scaffold box (yellow)
\newtcolorbox{scaffold}[1]{
  breakable, arc=2.5mm, boxrule=0.5pt,
  colback=GoldLight, colframe=Amber!40,
  fonttitle=\headingfont\bfseries, coltitle=GoldDark,
  title={#1},
  left=4mm, right=3.5mm, top=2.5mm, bottom=2.5mm
}

% Worked example box
\newtcolorbox{workedexbox}[1]{
  breakable, arc=2.5mm, boxrule=0.5pt,
  colback=white, colframe=Amber!50,
  fonttitle=\headingfont\bfseries, coltitle=GoldDark,
  title={#1},
  left=4mm, right=3.5mm, top=2.5mm, bottom=2.5mm
}
% Worked examples: \begin{workedex}{title} -> "Worked Example N.M --- title",
% auto-numbered per section and labelled we:N.M. Never type the number.
\newcounter{wenum}[section]
\renewcommand{\thewenum}{\thesection.\arabic{wenum}}
\newenvironment{workedex}[1]
  {\refstepcounter{wenum}\label{we:\thewenum}\begin{workedexbox}{Worked Example \thewenum\ --- #1}}
  {\end{workedexbox}}

% We do: straight after each worked example, a near-identical question the
% class works through together, with its own writing space. Numbered to
% match the worked example above it and labelled wd:N.M.
%   \begin{wedo}{title} question ... \answerbox{h} \end{wedo}
\newtcolorbox{wedobox}[1]{
  enhanced, breakable=false, arc=2.5mm, boxrule=0.5pt,
  colback=white, colframe=TealGreen!55,
  fonttitle=\headingfont\bfseries, coltitle=TealGreen!75!black,
  title={#1},
  left=4mm, right=3.5mm, top=2.5mm, bottom=2.5mm
}
\newenvironment{wedo}[1]
  {\par\needspace{9\baselineskip}\begin{wedobox}{We do \thewenum\ --- #1}\label{wd:\thewenum}\ignorespaces}
  {\end{wedobox}}

% Year 7 box (blue)
\newtcolorbox{yr7box}{
  breakable, arc=2.5mm, boxrule=0.5pt,
  colback=LightBlue, colframe=Blue!40,
  left=4mm, right=3.5mm, top=2.5mm, bottom=2.5mm
}

% Year 8 extension box (violet)
\newtcolorbox{yr8box}{
  breakable, arc=2.5mm, boxrule=0.5pt,
  colback=PurpleLight, colframe=Violet!40,
  left=4mm, right=3.5mm, top=2.5mm, bottom=2.5mm
}

% Answer lines box
\newcommand{\answerbox}[1]{%
  \begin{tcolorbox}[
    enhanced, breakable, arc=2mm,
    colback=Paper, colframe=Line,
    left=2mm, right=2mm, top=1mm, bottom=1mm,
    boxrule=1pt,
    height=#1
  ]
  \end{tcolorbox}%
}

% Experiment header
\newtcolorbox{expheader}[3]{
  breakable=false, arc=2.5mm,
  colback=RowAlt, colframe=Amber!50, boxrule=0.5pt,
  fonttitle=\headingfont\bfseries\large, coltitle=NavyBlue,
  title={Experiment #1: #2},
  left=4mm, right=3.5mm, top=2.5mm, bottom=2.5mm,
  after upper={\par\smallskip
    {\headingfont\bfseries\color{GoldDark}Aim:}\ \textit{#3}},
}

% Reading comprehension passage box (green)
\newtcolorbox{readingbox}[1]{
  breakable, arc=2.5mm, boxrule=0.5pt,
  colback=white, colframe=Teal!50,
  fonttitle=\headingfont\bfseries, coltitle=TealGreen,
  title={#1},
  left=4mm, right=3.5mm, top=2.5mm, bottom=2.5mm,
  fontupper=\normalsize
}

% Skill-builder box (for beginner scaffolding)
\newtcolorbox{skillbox}[1]{
  breakable, arc=2.5mm, boxrule=0.5pt,
  colback=RowAlt, colframe=Blue!40,
  fonttitle=\headingfont\bfseries, coltitle=MidBlue,
  title={#1},
  left=4mm, right=3.5mm, top=2.5mm, bottom=2.5mm
}

% ── Custom Commands ───────────────────────────────────────────────────────────

% Question: \question{marks}{text}  (marks can be empty)
% Numbered automatically as Q<section>.<n>, reset every section.
% Do NOT write numbers by hand.
\newcounter{qnum}[section]
\renewcommand{\theqnum}{\thesection.\arabic{qnum}}
\newcommand{\question}[2]{%
  \refstepcounter{qnum}%
  \needspace{3\baselineskip}%
  \medskip
  \noindent\label{q:\theqnum}\textbf{Q\theqnum}\ifx&#1&\else\ \textcolor{DarkGrey}{(#1 mark\ifnum#1=1\else s\fi)}\fi\quad #2%
  \medskip
}

% Inline marks note
\newcommand{\markscount}[1]{\hfill\textcolor{DarkGrey}{\textit{(#1 mark\ifnum#1=1\else s\fi)}}}

% ── Practice set counter and question macros ─────────────────────────────────
\newcounter{pqnum}
\newcommand{\practiceset}[1]{%
  \needspace{6\baselineskip}
  \setcounter{pqnum}{0}%
  \vspace{6pt}
  {\headingfont\bfseries\large\color{GoldDark} Practice Set --- #1}\\[-2pt]
  {\color{Amber}\rule{\linewidth}{1pt}}\\[2pt]
}
% \pq{question text}{working box height}
\newcommand{\pq}[2]{%
  \stepcounter{pqnum}%
  \needspace{4\baselineskip}
  \vspace{5pt}
  \noindent\textbf{\thepqnum.}\quad #1\par
  \vspace{2pt}
  \answerbox{#2}
}
% Two short questions side by side, each with its own box
\newcommand{\pqpair}[4]{%
  \needspace{4\baselineskip}
  \vspace{5pt}
  \noindent
  \begin{minipage}[t]{0.48\linewidth}
    \stepcounter{pqnum}\textbf{\thepqnum.}\quad #1\par\vspace{2pt}
    \begin{tcolorbox}[enhanced,arc=2mm,colback=Paper,colframe=Line,left=2mm,right=2mm,
      top=1mm,bottom=1mm,boxrule=1pt,height=#2]\end{tcolorbox}
  \end{minipage}\hfill
  \begin{minipage}[t]{0.48\linewidth}
    \stepcounter{pqnum}\textbf{\thepqnum.}\quad #3\par\vspace{2pt}
    \begin{tcolorbox}[enhanced,arc=2mm,colback=Paper,colframe=Line,left=2mm,right=2mm,
      top=1mm,bottom=1mm,boxrule=1pt,height=#4]\end{tcolorbox}
  \end{minipage}
}

% ── Comprehension question macro (compact, lined) ────────────────────────────
\newcounter{rcnum}
\newcommand{\rcreset}{\setcounter{rcnum}{0}}
\newcommand{\rcq}[2]{%
  \stepcounter{rcnum}%
  \needspace{3\baselineskip}
  \vspace{4pt}
  \noindent\textbf{C\thercnum.}\quad #1\par
  \vspace{2pt}
  \linedbox{#2}
}

% Section divider (big numbered banner): \sectiondiv{title}{summary}
%
% This IS the section: it steps the real section counter, so subsections,
% tasks, questions and figures inside it are numbered N.M automatically.
% Put \label{sec:N} straight after it. Never hand-type a section number.
\newcommand{\sectiondiv}[2]{%
  \needspace{5\baselineskip}%
  \refstepcounter{section}%
  % PDF bookmark for the lesson; its anchor, lesson.N.1, is what the contents table links to
  \pdfbookmark[1]{Lesson \thesection{} — #1}{lesson.\thesection}%
  \begin{tcolorbox}[
    enhanced, breakable=false, arc=3.5mm,
    colback=NavyBlue, colframe=NavyBlue, boxrule=0pt,
    sidebyside, sidebyside align=center,
    lefthand width=2.4cm,
    left=2mm, right=4mm, top=4mm, bottom=4mm
  ]
    \begin{center}
      {\headingfont\bfseries\color{white!35!NavyBlue}\fontsize{50}{50}\selectfont \thesection}
    \end{center}
    \tcblower
    {\headingfont\bfseries\color{white}\LARGE #1}\\[3pt]
    {\itshape\color{white!75!NavyBlue}\small #2}
  \end{tcolorbox}%
}

% Tasks: \task, \task[a] (part a of a NEW task), \taskpart{b} (next part
% of the same task), \exttask (Year 8 "Extension" — same counter).
% Numbered Task <section>.<n>, reset every section. Do NOT type numbers.
% Every task, question and figure also gets an automatic label
% (task:1.7b, q:2.4, fig:1.2, fig:E2.1) so page refs come from the .aux.
\newcounter{tasknum}[section]
\renewcommand{\thetasknum}{\thesection.\arabic{tasknum}}
\NewDocumentCommand{\task}{o}{\refstepcounter{tasknum}\label{task:\thetasknum\IfValueT{#1}{#1}}\textbf{Task \thetasknum\IfValueT{#1}{#1}}}
\newcommand{\taskpart}[1]{\label{task:\thetasknum#1}\textbf{Task \thetasknum#1}}
\newcommand{\exttask}{\refstepcounter{tasknum}\label{task:\thetasknum}\textbf{Extension \thetasknum}}

% Figure captions: \figcap[fig:label]{caption} -> "Figure N.M --- caption"
% \expfigcap{caption} -> "Figure EN.M" (experiment pages). Both auto-numbered.
\newcounter{fignum}[section]
\renewcommand{\thefignum}{\thesection.\arabic{fignum}}
\newcounter{expfignum}[section]
\renewcommand{\theexpfignum}{E\thesection.\arabic{expfignum}}
\newcommand{\figcap}[2][]{%
  \begin{center}
    \refstepcounter{fignum}\label{fig:\thefignum}\if\relax\detokenize{#1}\relax\else\label{#1}\fi
    \small\itshape\color{DarkGrey} Figure \thefignum\ --- #2
  \end{center}%
}
\newcommand{\expfigcap}[1]{%
  \begin{center}
    \refstepcounter{expfignum}\label{fig:\theexpfignum}%
    \small\itshape\color{DarkGrey} Figure \theexpfignum\ --- #1
  \end{center}%
}

% Diagram include
\newcommand{\diagramfig}[3]{%  {filename}{width fraction}{caption}
  \begin{center}
    \includegraphics[width=#2\linewidth]{#1}\\[4pt]
    \small\itshape\color{DarkGrey} #3
  \end{center}%
}

% Table column types
\newcolumntype{L}[1]{>{\RaggedRight\arraybackslash}p{#1}}
\newcolumntype{C}[1]{>{\Centering\arraybackslash}p{#1}}

% Table header colour
\newcommand{\thead}[1]{\cellcolor{TableHead}\color{white}\bfseries #1}
\newcommand{\talt}{\rowcolor{RowAlt}}
\newcommand{\tmid}{\rowcolor{MidGrey}}

% Inline year tags
\newcommand{\tagyr}[1]{%
  \tikz[baseline=-0.5ex]\node[
    rounded corners=2pt,
    fill=#1!20, draw=#1, text=#1,
    inner xsep=4pt, inner ysep=1pt,
    font=\scriptsize\bfseries
  ]{Yr #1 only};}

% Draw answer lines
\newcommand{\anslines}[1]{%
  \vspace{2pt}%
  \foreach \n in {1,...,#1}{%
    \par\noindent\rule{\linewidth}{0.3pt}\vspace{3pt}%
  }%
}

% ── Document ─────────────────────────────────────────────────────────────────
% ── Optional-lesson note, shown under the lesson banner ──────────────────────
\newcommand{\optionalnote}[1]{%
  % the title is a bold lead-in on the first line rather than a title bar,
  % which keeps Lesson 11's reading on the lesson's first page
  \begin{tcolorbox}[enhanced, breakable=false, arc=2.5mm, boxrule=0.5pt,
    colback=PurpleLight, colframe=Violet!50,
    left=4mm, right=3.5mm, top=2mm, bottom=2mm]
  {\headingfont\bfseries\color{PurpleDark}Optional lesson \textbullet{} Year 8 extension.}\enspace #1
  \end{tcolorbox}\vspace{2pt}}


% ══════════════════════════════════════════════════════════════════════════════
%  LESSON STRUCTURE (lesson booklet, Oct 2026)
%  Every lesson runs in the same four parts, in this order:
%    \lessonpart{A}{Reading and comprehension}{Teal}   \label{read:N}
%    \lessonpart{B}{Theory and examples}{Blue}         \label{theory:N}
%    \lessonpart{C}{Practice}{Amber}                    \label{prac:N}
%    \lessonpart{D}{Above level}{Violet}                \label{ext:N}
%  Above-level / higher-order questions live ONLY in Part D, at the end.
% ══════════════════════════════════════════════════════════════════════════════
\newcommand{\lessonpart}[3]{%
  \par\if#1A\needspace{10\baselineskip}\vspace{10pt}\else
    \if#1C\notesfill\fi\clearpage\fi
  \begin{tcolorbox}[enhanced, breakable=false, arc=2mm, boxrule=0pt,
    colback=#3!14, borderline west={4pt}{0pt}{#3},
    left=4mm, right=3mm, top=1.6mm, bottom=1.6mm]
  {\headingfont\bfseries\color{#3!60!black}\large Part #1 \textbullet{} #2}
  \end{tcolorbox}\vspace{2pt}}

% Key vocabulary at the start of each lesson: students copy each definition
% from the vocabulary slide in the lesson's starter. First row = red scaffold.
\newcommand{\vocabhead}{%
  \par\vspace{6pt}%
  \begin{tcolorbox}[enhanced, breakable=false, arc=2mm, boxrule=0pt,
    colback=NavyBlue!10, borderline west={4pt}{0pt}{NavyBlue},
    left=4mm, right=3mm, top=1.6mm, bottom=1.6mm]
  {\headingfont\bfseries\color{NavyBlue}\large Key vocabulary}
  \end{tcolorbox}\vspace{2pt}%
  \noindent Copy the definition of each word from the slide. The first one has been done for you.\par}

% Curriculum codes under a lesson banner
\newcommand{\lessoncodes}[1]{%
  {\small\headingfont\color{DarkGrey}Victorian Curriculum 2.0:\ \textbf{#1}}\par\vspace{4pt}}

% Two short tasks side by side, each with its own working box
% \taskpair{text}{box height}{text}{box height}
% The two boxes line up: the shorter question is padded by the difference in
% text height (measured in a scratch box, task counter restored, no label).
\newsavebox\tpbox
\newlength\tpA \newlength\tpB
\newcount\tpsaved
\newcommand\tpmeasure[2]{%
  \tpsaved=\value{tasknum}%
  \sbox\tpbox{\vbox{\hsize=0.485\linewidth \linewidth=\hsize \parskip=0pt
    \renewcommand\label[1]{}\task #1\par}}%
  \setlength{#2}{\dimexpr\ht\tpbox+\dp\tpbox\relax}%
  \setcounter{tasknum}{\tpsaved}}
\newcommand{\taskpair}[4]{%
  \par\needspace{9\baselineskip}\vspace{4pt}%
  \tpmeasure{#1}{\tpA}\tpmeasure{#3}{\tpB}%
  \noindent
  \edef\abgrp{\number\numexpr\value{abox}+1\relax}%
  \begin{minipage}[t]{0.485\linewidth}
    \task #1\par\ifdim\tpB>\tpA\vspace{\dimexpr\tpB-\tpA\relax}\fi\vspace{2pt}\answerbox{#2}
  \end{minipage}\hfill
  \begin{minipage}[t]{0.485\linewidth}
    \task #3\par\ifdim\tpA>\tpB\vspace{\dimexpr\tpA-\tpB\relax}\fi\vspace{2pt}\answerbox{#4}
  \end{minipage}\par\gdef\abgrp{0}\vspace{4pt}}
% Two parts of one question side by side, with their answer boxes level:
% \ptpair{a}{part (a) text}{box height}{b}{part (b) text}{box height}.
% The shorter part's text gets space below it to match the longer one.
\newcommand\ptmeasure[3]{%
  \sbox\tpbox{\vbox{\hsize=0.485\linewidth \linewidth=\hsize \parskip=0pt \pt{#1}#2\par}}%
  \setlength{#3}{\dimexpr\ht\tpbox+\dp\tpbox\relax}}
\newcommand{\ptpair}[6]{%
  \ptmeasure{#1}{#2}{\tpA}\ptmeasure{#4}{#5}{\tpB}%
  \par\noindent
  \edef\abgrp{\number\numexpr\value{abox}+1\relax}%
  \begin{minipage}[t]{0.485\linewidth}
    \pt{#1}#2\par\ifdim\tpB>\tpA\vspace{\dimexpr\tpB-\tpA\relax}\fi\answerbox{#3}
  \end{minipage}\hfill
  \begin{minipage}[t]{0.485\linewidth}
    \pt{#4}#5\par\ifdim\tpA>\tpB\vspace{\dimexpr\tpA-\tpB\relax}\fi\answerbox{#6}
  \end{minipage}\par\gdef\abgrp{0}}

% Answer boxes never split across a page (a half box on its own page is useless)
% ── Page filling ────────────────────────────────────────────────────────────
% Every writing box gets an id. Its page is logged to \jobname.abp; fill.py
% measures the empty space left at the foot of each page and shares it out
% among that page's boxes via \jobname.abx (\abset{id}{extra}). Side-by-side
% boxes share a group so they grow together.
\newcounter{abox}
\newwrite\abpos \immediate\openout\abpos=\jobname.abp
\newcommand{\abset}[2]{\expandafter\gdef\csname abx@#1\endcsname{#2}}
\InputIfFileExists{\jobname.abx}{}{}
\def\abgrp{0}
\newcommand{\abnext}{\stepcounter{abox}%
  \edef\abcur{\ifcsname abx@\theabox\endcsname\csname abx@\theabox\endcsname\else 0pt\fi}}
\newcommand{\ablog}{\edef\abtmp{\noexpand\write\abpos{\theabox\space\abgrp\space\noexpand\thepage}}\abtmp}
\newcommand{\abh}[1]{\dimexpr#1+\abcur\relax}
\renewcommand{\answerbox}[1]{%
  \abnext
  \begin{tcolorbox}[enhanced, breakable=false, arc=2mm, colback=Paper, colframe=Line,
    left=2mm, right=2mm, top=1mm, bottom=1mm, boxrule=1pt, height=\abh{#1}]
  \end{tcolorbox}\ablog%
}
% \linedbox{height}: an answer box ruled in dotted thirds. Lines are \thirdrow/3 apart and run dotted, dotted,
% solid (the baseline), repeating, and end on one dotted line. The number of
% rows is worked out from the box's final height, so the ruling still fills the
% box after fill.py has grown it; the rows are centred in the box.
\newlength\thirdrow \setlength\thirdrow{9mm}
\tikzset{thirdsolid/.style={draw=DarkGrey!70, line width=0.5pt},
         thirddot/.style={draw=DarkGrey!70, line width=0.6pt, dash pattern=on 0.9pt off 1.6pt}}
\newlength\lbH \newlength\lbO \newcount\lbN
\newcommand{\linedbox}[1]{%
  \abnext
  % interior height = box height less the 1pt frame top and bottom
  \setlength\lbH{\dimexpr\abh{#1}-2pt\relax}%
  % whole rows that fit with at least 3 mm spare top and bottom (at least one)
  \lbN=\numexpr(\dimexpr\lbH-6mm-\thirdrow/2\relax)/\thirdrow\relax
  \ifnum\lbN<1 \lbN=1 \fi
  \setlength\lbO{\dimexpr(\lbH-\lbN\thirdrow)/2\relax}%
  \begin{tcolorbox}[enhanced, breakable=false, arc=2mm, colback=Paper, colframe=Line,
    left=2mm, right=2mm, top=1mm, bottom=1mm, boxrule=1pt, height=\abh{#1},
    underlay={\begin{tcbclipinterior}
      \foreach \i in {1,...,\the\lbN}{
        \draw[thirddot] ([xshift=3mm,yshift=-\lbO-(\i-1)*\thirdrow]interior.north west) -- ([xshift=-3mm,yshift=-\lbO-(\i-1)*\thirdrow]interior.north east);
        \draw[thirddot] ([xshift=3mm,yshift=-\lbO-(\i-1)*\thirdrow-\thirdrow/3]interior.north west) -- ([xshift=-3mm,yshift=-\lbO-(\i-1)*\thirdrow-\thirdrow/3]interior.north east);
        \draw[thirdsolid] ([xshift=3mm,yshift=-\lbO-(\i-1)*\thirdrow-2*\thirdrow/3]interior.north west) -- ([xshift=-3mm,yshift=-\lbO-(\i-1)*\thirdrow-2*\thirdrow/3]interior.north east);}
      \draw[thirddot] ([xshift=3mm,yshift=-\lbO-\lbN\thirdrow]interior.north west) -- ([xshift=-3mm,yshift=-\lbO-\lbN\thirdrow]interior.north east);
    \end{tcbclipinterior}}]
  \end{tcolorbox}\ablog%
}
% Every written-response box is ruled the same way (Part A comprehension, We do,
% Part C tasks and Part D extensions, including side-by-side pairs). Space for
% mathematical working uses \workbox{height} instead and stays plain, as do
% \drawbox and the write-in tables.
\let\workbox\answerbox   % plain box: mathematical working (calculations) stays unruled
\let\answerbox\linedbox
% \drawbox{height}{contents}: white box for a diagram or drawing
\newcommand{\drawbox}[2]{%
  \abnext
  \begin{tcolorbox}[colback=white, colframe=MidGrey, boxrule=0.5pt, height=\abh{#1}]#2\end{tcolorbox}\ablog}
% Notes box at the end of Part B: fills whatever is left of the page (if worth it)
\newcommand{\notesfill}{\par
  \ifdim\dimexpr\pagegoal-\pagetotal\relax>3.2cm
    \vspace*{\fill}%
    \begin{tcolorbox}[enhanced, breakable=false, arc=2mm, colback=Paper, colframe=Line,
      left=2mm, right=2mm, top=1mm, bottom=1mm, boxrule=1pt,
      height=\dimexpr\pagegoal-\pagetotal-1.0cm\relax]
      {\small\headingfont\color{DarkGrey}Notes}
    \end{tcolorbox}%
  \fi}



% Centred table: wrap every content table in \begin{ctab} ... \end{ctab}
% (figures already centre themselves inside their \diag... macros)
% ctab also registers write-in tables (S column) for page filling: the extra
% height is per row, logged with the row count.
\newenvironment{ctab}{\par\smallskip\begingroup\centering
  \abnext\let\tbcur\abcur\global\wtrows0 }{\par
  \ifnum\wtrows>0 \edef\abtmp{\noexpand\write\abpos{\theabox\space\abgrp\space\noexpand\thepage\space\the\wtrows}}\abtmp\fi
  \endgroup\smallskip}

% ── Booklet formatting rules (Oct 2026) ─────────────────────────────────────
% 1. No marks are shown anywhere in the booklet.
\renewcommand{\markscount}[1]{}
% 2. A task's heading sits on its own line; its text starts on the next
%    line, left aligned. Parts (a), (b)… each start a new line: use \pt{a}.
\RenewDocumentCommand{\task}{o}{\par\noindent\refstepcounter{tasknum}\label{task:\thetasknum\IfValueT{#1}{#1}}\textbf{Task \thetasknum\IfValueT{#1}{#1}}\par\nobreak\noindent\ignorespaces}
\renewcommand{\taskpart}[1]{\par\noindent\label{task:\thetasknum#1}\textbf{Task \thetasknum#1}\par\nobreak\noindent\ignorespaces}
\renewcommand{\exttask}{\par\noindent\refstepcounter{tasknum}\label{task:\thetasknum}\textbf{Extension \thetasknum}\par\nobreak\noindent\ignorespaces}
\newcommand{\pt}[1]{\par\noindent\textbf{(#1)}\ \ignorespaces}
% Workstation pages (Experiment 1): numbered method steps with a hanging
% indent, and a one-line safety reminder under the materials list.
\newcommand{\wsstep}[1]{\par\noindent\hangindent=1.6em\hangafter=1\makebox[1.6em][l]{\textbf{#1}}\ignorespaces}
\newcommand{\wssafety}[1]{\par\smallskip{\small\textcolor{ScaffoldRed}{\headingfont\bfseries Safety:}\ #1}\par}
% 3. Tables have vertical rules between columns and down both sides
%    (write | between and around the columns, e.g. {|L{3cm}|C{2cm}|}).
%    Keep the normal padding at the outer edges (no @{} there): colortbl paints
%    header colour one \tabcolsep past each edge, so trimming the padding makes
%    the header overhang the grid.
\arrayrulecolor{MidGrey}
\setlength{\arrayrulewidth}{0.6pt}
% 4. Write-in tables: W/Z are left/centred columns that centre text
%    vertically; put the S column first (see below).
\newcolumntype{W}[1]{>{\RaggedRight\arraybackslash}m{#1}}
\newcolumntype{Z}[1]{>{\Centering\arraybackslash}m{#1}}
\def\tbcur{0pt}
\newcount\wtrows
\newcommand{\wh}{\global\advance\wtrows1 \rule[\dimexpr-0.75cm-\tbcur/2\relax]{0pt}{\dimexpr1.5cm+\tbcur\relax}}
% S = zero-width first column that makes every row of a write-in table
% at least 1.5 cm tall with its text vertically centred. Start each data
% row with "&"; start the header row with "\hd &" so it stays slim.
\newif\ifnowh
\newcolumntype{S}{@{}>{}c<{\ifnowh\else\wh\fi\global\nowhfalse}@{\hspace{\tabcolsep}}}
\newcommand{\hd}{\global\nowhtrue}
% 5. The first row of a multi-row write-in table is a worked scaffold in red.
\definecolor{ScaffoldRed}{HTML}{C0392B}
\newcommand{\scaf}[1]{\textcolor{ScaffoldRed}{\itshape #1}}
% 6. Tables are a light grid (rules between rows and columns) on white —
%    row shading is off, because colortbl paints shading over column rules.
\renewcommand{\talt}{}
\renewcommand{\tmid}{}
% ══════════════════════════════════════════════════════════════════════════════
%  EXPERIMENT BOOKLET LAYER (Nov 2026)
%  From 2027 every Year 7 unit collects its experiments into one A4 experiment
%  booklet that puts the experimental techniques and skills first. Every skill
%  comes from the faculty Science Skills continuum (Years 7--8 column) and is
%  named the same way everywhere: a code (QP1, PC4, PA2 …) on a chip coloured by
%  its sub-strand. Copy this block into the next unit's experiment booklet.
%
%  Each experiment runs in the same order:
%    \experimentdiv{title}{summary}{lesson}\label{sec:N}
%    key vocabulary · learning intentions · \begin{skillsfocus}
%    \exppart*{Technique spotlight}{Blue}   (technique, worked example, We do)
%    \exppart{Plan and safety}{Teal}        (aim, prediction, variables, risks)
%    \exppart{Method}{Teal}
%    \exppart{Results}{Amber}
%    \exppart{Analysis and evaluation}{Coral}
%    \exppart{Skill check}{NavyBlue}
% ══════════════════════════════════════════════════════════════════════════════
\usetikzlibrary{arrows.meta,calc,positioning}
\usepackage{pifont}

\newenvironment{keepext}{\par\begin{tcolorbox}[blankest, breakable=false]}{\end{tcolorbox}\par}
% graph-paper box: \graphgrid{height}
\newcommand{\graphgrid}[1]{%
  \begin{tcolorbox}[enhanced, breakable=false, colback=white, colframe=MidGrey, boxrule=0.5pt, height=#1,
    underlay={\begin{tcbclipinterior}\draw[step=5mm, Line, very thin] (interior.south west) grid (interior.north east);
      \draw[step=1cm, MidGrey!70, thin] (interior.south west) grid (interior.north east);\end{tcbclipinterior}}]
  \end{tcolorbox}}

% ── Skills (faculty Science Skills continuum, Years 7–8) ─────────────────────
% \defskill{code}{name}{colour}{VC code}{what it means}
\makeatletter
\newcommand{\defskill}[5]{%
  \expandafter\gdef\csname sk@n@#1\endcsname{#2}%
  \expandafter\gdef\csname sk@c@#1\endcsname{#3}%
  \expandafter\gdef\csname sk@v@#1\endcsname{#4}%
  \expandafter\gdef\csname sk@d@#1\endcsname{#5}}
\newcommand{\skname}[1]{\csname sk@n@#1\endcsname}
\newcommand{\skcode}[1]{\csname sk@v@#1\endcsname}
\newcommand{\skdesc}[1]{\csname sk@d@#1\endcsname}
% \sk{PC4}: the chip — code and name on the sub-strand colour
\newcommand{\sk}[1]{%
  \edef\sk@col{\csname sk@c@#1\endcsname}%
  \tikz[baseline=(c.base)]\node[rounded corners=2.5pt, fill=\sk@col!16, draw=\sk@col, line width=0.6pt,
    text=\sk@col!55!black, inner xsep=4pt, inner ysep=2pt, font=\headingfont\bfseries\scriptsize]
    (c) {#1\,\textperiodcentered\,\skname{#1}};}
% \skr{PC4}: the chip flush right at the end of a paragraph; drops to its own line if it does not fit
\newcommand{\skr}[1]{\unskip\nobreak\hfill\penalty50\hskip1em\null\nobreak\hfill\sk{#1}}
% \skc{PC4}: the code alone, for tight table cells
\newcommand{\skc}[1]{%
  \edef\sk@col{\csname sk@c@#1\endcsname}%
  \tikz[baseline=(c.base)]\node[rounded corners=2.5pt, fill=\sk@col!16, draw=\sk@col, line width=0.6pt,
    text=\sk@col!55!black, inner xsep=3pt, inner ysep=2pt, font=\headingfont\bfseries\scriptsize]
    (c) {#1};}
\makeatother

% Questioning and predicting — VC2S8I01
\defskill{QP1}{Testable hypothesis}{Teal}{VC2S8I01}{A prediction an experiment can test, naming both variables and the direction of change: \emph{If the [IV] is \dots, the [DV] will \dots}; a reason is optional}
\defskill{QP2}{Formulate an aim}{Teal}{VC2S8I01}{One sentence naming both variables: \emph{To investigate the effect of changing [IV] on [DV]}}
% Planning and conducting — VC2S8I02, VC2S8I03
\defskill{PC1}{Plan and follow a method}{Blue}{VC2S8I02}{Follow, and later write, numbered steps in order: named equipment, quantities with units, no personal pronouns and a labelled diagram, so that someone else could repeat it (repeatable)}
\defskill{PC2}{Identify variables}{Blue}{VC2S8I02}{Name the independent variable (what you change), the dependent variable (what you measure) and at least two controlled variables (what you keep the same)}
\defskill{PC3}{Risks and safety}{Blue}{VC2S8I02}{Name each risk, how to prevent it and how to respond if it happens; follow the lab rules}
\defskill{PC4}{Precision}{Blue}{VC2S8I02}{Choose equipment with a fine enough scale, read it to its smallest division, and repeat measurements to see how closely they agree}
\defskill{PC5}{Ethics}{Blue}{VC2S8I02}{Treat people, living things, the environment and other people's data with respect}
\defskill{PC6}{Collect data and measure}{Blue}{VC2S8I02}{Take each measurement at the right time, with its unit, and record it straight away}
\defskill{PC7}{Use and select equipment}{Blue}{VC2S8I03}{Name each piece of equipment, use it correctly and choose the right piece for the job}
\defskill{PC8}{Use digital tools}{Blue}{VC2S8I03}{Use stopwatches, data loggers, sensors and spreadsheets to collect and record data}
% Processing, modelling and analysing — VC2S8I04, VC2S8I05
\defskill{PA1}{Tables}{GoldDark}{VC2S8I04}{Title, Headings, Units in the headings only, Decimals kept the same (THUD), with the independent variable in the first column}
\defskill{PA2}{Graphing}{GoldDark}{VC2S8I04}{The right type of graph, a title, axes labelled with units, the independent variable on the x-axis, an even scale and accurate points, then join the points or draw a line of best fit}
\defskill{PA3}{Mathematical relationships}{GoldDark}{VC2S8I04}{Process data with calculations: differences, averages, rates and formulas}
\defskill{PA4}{Creating models}{GoldDark}{VC2S8I04}{Build a physical, drawn or mathematical model that represents something real}
\defskill{PA5}{Patterns and trends}{GoldDark}{VC2S8I05}{Describe how the dependent variable changes as the independent variable changes, quoting data}
\defskill{PA6}{Correlation and causality}{GoldDark}{VC2S8I05}{Say whether two variables change together, and whether there is evidence that one causes the other}
\defskill{PA7}{Identify anomalies}{GoldDark}{VC2S8I05}{Spot a result that does not fit the pattern and suggest why it happened}
\defskill{PA8}{Interpreting models}{GoldDark}{VC2S8I05}{Use a model to explain an observation, and say what the model shows well and what it leaves out}
% Evaluating — VC2S8I06, VC2S8I07
\defskill{EV1}{Results against hypothesis}{OrangeDark}{VC2S8I06}{State the general trend and whether it supports the hypothesis (supported, never proven), with data to back the trend}
\defskill{EV2}{Compare with known concepts}{OrangeDark}{VC2S8I06}{Compare your results with accepted values and scientific ideas}
\defskill{EV3}{Error sources}{OrangeDark}{VC2S8I06}{Name what could have made results less accurate or less precise, and suggest how to improve the method}
\defskill{EV4}{Quality of secondary sources}{OrangeDark}{VC2S8I07}{Judge whether information from a website, book or video can be trusted}
\defskill{EV5}{Questions for further investigation}{OrangeDark}{VC2S8I07}{Pose a new question that follows from your results}
% Communicating — VC2S8I08
\defskill{CO1}{Scientific writing}{PurpleDark}{VC2S8I08}{Write in the third person, legibly, with correct spelling and grammar, short and to the point, using scientific terms}
\defskill{CO2}{Presentation formats}{PurpleDark}{VC2S8I08}{Choose the right format to show findings: a table, graph, scientific diagram, poster, report or talk}

% Skills focus panel at the top of each experiment:
%   \begin{skillsfocus} \focus{PC4}{what you do with it here} … \end{skillsfocus}
\newtcolorbox{skillsfocus}{
  enhanced, breakable=false, arc=2.5mm, boxrule=0.6pt,
  colback=white, colframe=NavyBlue!55,
  fonttitle=\headingfont\bfseries, coltitle=white, colbacktitle=NavyBlue,
  title={Skills focus --- the science skills this experiment builds},
  left=4mm, right=3.5mm, top=2mm, bottom=2mm}
\newcommand{\focus}[2]{\par\noindent\makebox[1.25cm][l]{\skc{#1}}\parbox[t]{\dimexpr\linewidth-1.25cm\relax}{\small\textbf{\skname{#1}}\,---\,#2}\par\vspace{2pt}}

% Learning intentions panel: \begin{lipanel} \inquiryli{…} \understandli{…} \sclist{…} \end{lipanel}
\newtcolorbox{lipanel}{
  enhanced, breakable=false, arc=2.5mm, boxrule=0.5pt,
  colback=LightBlue, colframe=Blue!40,
  fonttitle=\headingfont\bfseries, coltitle=MidBlue, title={Learning intentions},
  left=4mm, right=3.5mm, top=2mm, bottom=2mm}
\newcommand{\inquiryli}[1]{\par\noindent{\headingfont\bfseries\color{NavyBlue}Science Inquiry:}\ #1\par\vspace{2pt}}
\newcommand{\understandli}[1]{\par\noindent{\headingfont\bfseries\color{NavyBlue}Science Understanding:}\ #1\par\vspace{2pt}}
\newcommand{\sctitle}{\par\noindent{\headingfont\bfseries\color{NavyBlue}Success criteria}\par}

% Technique spotlight: \begin{technique}{title}{PC4} … \end{technique}
\newtcolorbox{technique}[2]{
  enhanced, breakable, arc=2.5mm, boxrule=0.6pt,
  colback=white, colframe=Blue!55,
  fonttitle=\headingfont\bfseries, coltitle=MidBlue, colbacktitle=LightBlue,
  title={Technique: #1\hfill\sk{#2}},
  left=4mm, right=3.5mm, top=2.5mm, bottom=2.5mm}

% Experiment divider: \experimentdiv{title}{summary}{lesson number}
% Steps the section counter, so tasks, worked examples and figures are N.M.
\newcommand{\experimentdiv}[3]{%
  \clearpage
  \refstepcounter{section}%
  \pdfbookmark[1]{Experiment \thesection{} — #1}{exp.\thesection}%
  \begin{tcolorbox}[
    enhanced, breakable=false, arc=3.5mm,
    colback=NavyBlue, colframe=NavyBlue, boxrule=0pt,
    sidebyside, sidebyside align=center,
    lefthand width=2.6cm,
    left=2mm, right=4mm, top=4mm, bottom=4mm]
    \begin{center}
      {\headingfont\bfseries\color{white!60!NavyBlue}\scriptsize EXPERIMENT}\\[-2pt]
      {\headingfont\bfseries\color{white!35!NavyBlue}\fontsize{46}{46}\selectfont \thesection}
    \end{center}
    \tcblower
    {\headingfont\bfseries\color{white}\LARGE #1}\\[3pt]
    {\itshape\color{white!75!NavyBlue}\small #2}\\[3pt]
    {\headingfont\color{Sky}\footnotesize Lesson #3}
  \end{tcolorbox}%
}

% Part bar inside an experiment. \exppart{title}{colour} starts a new page
% (new parts always start on a new page); \exppart*{…}{…} follows straight on.
\NewDocumentCommand{\exppart}{s m m}{%
  \par\IfBooleanTF{#1}{\needspace{10\baselineskip}\vspace{8pt}}{\clearpage}%
  \begin{tcolorbox}[enhanced, breakable=false, arc=2mm, boxrule=0pt,
    colback=#3!14, borderline west={4pt}{0pt}{#3},
    left=4mm, right=3mm, top=1.6mm, bottom=1.6mm]
  {\headingfont\bfseries\color{#3!60!black}\large #2}
  \end{tcolorbox}\vspace{2pt}}

% Skill check at the end of each experiment: students tick one column per row.
%   \begin{skillcheck} \scrow{PC4}{I can …} … \end{skillcheck}
\newcommand{\tickbox}{\tikz[baseline=-0.6ex]\draw[DarkGrey, line width=0.7pt, rounded corners=1pt] (0,-0.17) rectangle (0.34,0.17);}
\newenvironment{skillcheck}{%
  \par\noindent Read each statement. Tick the column that best describes you \emph{today}.\par
  \begin{ctab}\renewcommand{\arraystretch}{1.25}%
  \begin{tabular}{|W{3.9cm}|W{7.0cm}|Z{1.45cm}|Z{1.45cm}|Z{1.45cm}|}
  \hline\thead{Skill} & \thead{Success criterion} & \thead{Not yet} & \thead{Nearly} & \thead{Got it} \\ \hline}
  {\end{tabular}\end{ctab}}
\newcommand{\scrow}[2]{\rule{0pt}{1.1cm}\skc{#1}\enspace{\small\skname{#1}} & #2 & \tickbox & \tickbox & \tickbox \\ \hline}

% Write-in label under a drawing: \nameline{label}
\newcommand{\nameline}[1]{\par\noindent{\small #1}\ \leaders\hbox{\rule[-2pt]{1pt}{0.4pt}}\hfill\null\par}

% ── Equipment drawn as scientific diagrams (2D, side-on, outline only) ──────
% Each \eq... macro draws at the origin inside a tikzpicture. \eqpic[scale]{\eqBeaker}
% puts one on its own. The same shapes are used everywhere, so the booklet
% models exactly the drawings students are asked to make.
\tikzset{
  eq/.style={line width=0.9pt, draw=Ink, line cap=round, line join=round},
  eqlab/.style={draw=Ink, line width=0.4pt},
  eqtext/.style={font=\footnotesize, text=Ink, inner sep=1.5pt},
}
\newcommand{\eqpic}[2][1]{\begin{tikzpicture}[scale=#1, baseline=(current bounding box.center)]#2\end{tikzpicture}}
\newcommand{\eqBeaker}{\draw[eq] (-0.22,2.62) -- (0,2.45) -- (0,0) -- (1.8,0) -- (1.8,2.5);}
\newcommand{\eqFlask}{\draw[eq] (0.62,2.9) -- (0.65,2.8) -- (0.65,1.9) -- (0,0) -- (1.9,0) -- (1.25,1.9) -- (1.25,2.8) -- (1.28,2.9);}
\newcommand{\eqTestTube}{\draw[eq] (-0.06,3.0) -- (0,2.9) -- (0,0.3) arc[start angle=180,end angle=360,radius=0.3] -- (0.6,2.9) -- (0.66,3.0);}
\newcommand{\eqCylinder}{%
  \draw[eq] (-0.15,3.6) -- (0,3.5) -- (0,0.15) (0.6,0.15) -- (0.6,3.5);
  \draw[eq] (-0.35,0) rectangle (0.95,0.15);
  \foreach \y in {0.6,1.1,...,3.2} \draw[Ink, line width=0.5pt] (0.6,\y) -- (0.4,\y);}
\newcommand{\eqFunnel}{\draw[eq] (0,2.2) -- (0.75,1.0) -- (0.75,0) (1.85,2.2) -- (1.1,1.0) -- (1.1,0);}
\newcommand{\eqBasin}{\draw[eq] (0,0.9) .. controls (0.15,0.02) and (2.25,0.02) .. (2.4,0.9);}
\newcommand{\eqTripod}{%
  \draw[eq] (-0.15,2.0) rectangle (2.55,2.12);
  \draw[eq] (0.25,2.0) -- (0,0) (2.15,2.0) -- (2.4,0) (1.2,2.0) -- (1.2,0.15);}
\newcommand{\eqBunsen}{%
  \draw[eq] (0,0) rectangle (1.6,0.22);
  \draw[eq] (0.58,0.85) -- (0.58,2.3) (1.02,0.85) -- (1.02,2.3);
  \draw[eq] (0.5,0.22) rectangle (1.1,0.85);
  \fill[Ink] (0.72,0.42) rectangle (0.88,0.62);
  \draw[eq] (-0.65,0.06) rectangle (0,0.16);}
\newcommand{\eqThermometer}{%
  \draw[eq] (0.25,0.25) circle (0.2);
  \draw[eq] (0.15,0.43) -- (0.15,3.9) arc[start angle=180,end angle=0,radius=0.1] -- (0.35,0.43);
  \draw[Ink, line width=1.3pt] (0.25,0.45) -- (0.25,1.9);}
\newcommand{\eqStand}{%
  \draw[eq] (0,0) rectangle (2.3,0.2);
  \draw[eq] (0.35,0.2) -- (0.35,4.0) (0.5,0.2) -- (0.5,4.0);
  \draw[eq] (0.25,2.75) rectangle (0.75,3.2);
  \draw[eq] (0.75,2.92) -- (1.6,2.92) (0.75,3.03) -- (1.6,3.03);
  \draw[eq] (1.6,2.7) -- (1.6,3.25) (1.6,2.7) -- (1.95,2.7) (1.6,3.25) -- (1.95,3.25);}
\newcommand{\eqWatchGlass}{\draw[eq] (0,0.35) .. controls (0.5,-0.05) and (1.7,-0.05) .. (2.2,0.35);}
\newcommand{\eqTongs}{%
  \draw[eq] (0,0.1) -- (2.0,0.45) .. controls (2.3,0.5) and (2.45,0.35) .. (2.6,0.15)
            (0,0.0) -- (2.0,-0.35) .. controls (2.3,-0.4) and (2.45,-0.25) .. (2.6,-0.05);}

% A label line from a point on a drawing to text at the side: \eqlab{(x,y)}{(tx,ty)}{text}
\newcommand{\eqlab}[3]{\draw[eqlab] #1 -- #2 node[eqtext, anchor=west] {#3};}
\newcommand{\eqlabl}[3]{\draw[eqlab] #1 -- #2 node[eqtext, anchor=east] {#3};}

% ── Reading scales ──────────────────────────────────────────────────────────
% \readcyl{lowest}{highest}{labelled step}{small divisions per label}{reading}
% A close-up of a measuring cylinder's scale, the water drawn to the reading:
% the bottom of the meniscus sits exactly on the reading.
\pgfkeys{/pgf/number format/.cd, fixed, precision=1}
\newcommand{\readcyl}[5]{%
\begin{tikzpicture}[baseline=(current bounding box.south)]
  \pgfmathsetmacro\rcH{5.0}
  \pgfmathsetmacro\rcY{(#5-#1)*\rcH/(#2-#1)}
  \pgfmathtruncatemacro\rcN{round((#2-#1)/#3*#4)}
  \fill[black!10] (0,-0.3) -- (0,\rcY+0.16) .. controls (0.45,\rcY-0.05) and (1.15,\rcY-0.05) .. (1.6,\rcY+0.16) -- (1.6,-0.3) -- cycle;
  \draw[Ink, line width=0.9pt] (0,\rcY+0.16) .. controls (0.45,\rcY-0.05) and (1.15,\rcY-0.05) .. (1.6,\rcY+0.16);
  \foreach \k in {0,...,\rcN}{
    \pgfmathsetmacro\yy{\k*\rcH/\rcN}
    \pgfmathtruncatemacro\ismaj{mod(\k,#4)==0 ? 1 : 0}
    \ifnum\ismaj=1
      \draw[Ink, line width=0.75pt] (0,\yy) -- (0.6,\yy);
      \pgfmathsetmacro\lv{#1+\k/#4*#3}
      \node[anchor=west, font=\scriptsize\bfseries, inner sep=1pt, fill=white, fill opacity=0.8, text opacity=1] at (0.66,\yy) {\pgfmathprintnumber{\lv}};
    \else
      \draw[Ink, line width=0.4pt] (0,\yy) -- (0.32,\yy);
    \fi}
  \draw[eq] (0,-0.3) -- (0,\rcH+0.35) (1.6,-0.3) -- (1.6,\rcH+0.35);
  \node[font=\scriptsize, anchor=north] at (0.8,-0.35) {scale in mL};
\end{tikzpicture}}
% \readtherm{lowest}{highest}{labelled step}{small divisions per label}{reading}
\newcommand{\readtherm}[5]{%
\begin{tikzpicture}[baseline=(current bounding box.south)]
  \pgfmathsetmacro\rtH{5.0}
  \pgfmathsetmacro\rtY{(#5-#1)*\rtH/(#2-#1)}
  \pgfmathtruncatemacro\rtN{round((#2-#1)/#3*#4)}
  \draw[eq, fill=white] (0,-0.45) circle (0.3);
  \fill[black!75] (0,-0.45) circle (0.22);
  \fill[black!75] (-0.07,-0.3) rectangle (0.07,\rtY);
  \draw[eq] (-0.16,-0.2) -- (-0.16,\rtH+0.3) arc[start angle=180,end angle=0,radius=0.16] -- (0.16,-0.2);
  \foreach \k in {0,...,\rtN}{
    \pgfmathsetmacro\yy{\k*\rtH/\rtN}
    \pgfmathtruncatemacro\ismaj{mod(\k,#4)==0 ? 1 : 0}
    \ifnum\ismaj=1
      \draw[Ink, line width=0.75pt] (0.25,\yy) -- (0.75,\yy);
      \pgfmathsetmacro\lv{#1+\k/#4*#3}
      \node[anchor=west, font=\scriptsize\bfseries, inner sep=1pt] at (0.8,\yy) {\pgfmathprintnumber{\lv}};
    \else
      \draw[Ink, line width=0.4pt] (0.25,\yy) -- (0.52,\yy);
    \fi}
  \node[font=\scriptsize, anchor=north] at (0.3,-0.8) {scale in \textdegree C};
\end{tikzpicture}}
% \balancefig{reading}: an electronic balance showing a reading
\newcommand{\balancefig}[1]{%
\begin{tikzpicture}[baseline=(current bounding box.south)]
  \draw[eq, rounded corners=3pt] (0,0) rectangle (4.2,1.35);
  \draw[eq] (1.9,1.35) -- (1.9,1.55) (2.3,1.35) -- (2.3,1.55);
  \draw[eq, rounded corners=1pt] (0.5,1.55) rectangle (3.7,1.72);
  \draw[eq, fill=LightGrey] (0.35,0.35) rectangle (2.75,1.02);
  \node[font=\headingfont\bfseries\large, anchor=east] at (2.65,0.68) {#1};
  \draw[eq] (3.3,0.68) circle (0.2);
  \node[font=\tiny\bfseries] at (3.3,0.3) {TARE};
\end{tikzpicture}}
% Bunsen burner with a flame: \bunsenflame{safety|heating}
\newcommand{\bunsenflame}[1]{%
\begin{tikzpicture}[baseline=(current bounding box.south)]
  \def\fl{#1}\def\safe{safety}%
  \ifx\fl\safe
    \fill[black!14] (0.8,2.3) .. controls (0.05,3.0) and (0.45,4.3) .. (0.85,5.0) .. controls (1.05,4.4) and (1.55,3.1) .. (0.8,2.3);
    \draw[Ink, line width=0.6pt] (0.8,2.3) .. controls (0.05,3.0) and (0.45,4.3) .. (0.85,5.0) .. controls (1.05,4.4) and (1.55,3.1) .. (0.8,2.3);
  \else
    \fill[white] (0.8,2.3) .. controls (0.45,2.8) and (0.6,3.5) .. (0.8,3.9) .. controls (1.0,3.5) and (1.15,2.8) .. (0.8,2.3);
    \draw[Ink, line width=0.6pt] (0.8,2.3) .. controls (0.45,2.8) and (0.6,3.5) .. (0.8,3.9) .. controls (1.0,3.5) and (1.15,2.8) .. (0.8,2.3);
    \fill[black!30] (0.8,2.3) .. controls (0.65,2.55) and (0.7,2.85) .. (0.8,3.0) .. controls (0.9,2.85) and (0.95,2.55) .. (0.8,2.3);
  \fi
  \draw[eq] (0,0) rectangle (1.6,0.22);
  \draw[eq] (0.58,0.85) -- (0.58,2.3) (1.02,0.85) -- (1.02,2.3);
  \draw[eq] (0.5,0.22) rectangle (1.1,0.85);
  \ifx\fl\safe\draw[Ink, line width=0.5pt, densely dotted] (0.72,0.42) rectangle (0.88,0.62);\else\fill[Ink] (0.72,0.42) rectangle (0.88,0.62);\fi
  \draw[eq] (-0.65,0.06) rectangle (0,0.16);
\end{tikzpicture}}
% ══════════════════════════════════════════════════════════════════════════════
%  BIODIVERSITY AND ECOSYSTEMS — COVER
%  The unit's cover art (covers/ecosystems.jpg), full bleed. The A5 lesson
%  booklet sets \wrcoverfile one folder deeper.
% ══════════════════════════════════════════════════════════════════════════════
\providecommand{\wrcoverfile}{../../covers/ecosystems.jpg}
\newcommand{\wrcover}{%
\begin{tikzpicture}[remember picture, overlay]
  \node[anchor=north] at (current page.north) {\includegraphics[width=\paperwidth]{\wrcoverfile}};
\end{tikzpicture}}
% ══════════════════════════════════════════════════════════════════════════════
%  2027 WORKBOOK STYLE (Year 7 workbook redesign, Nov 2026) — A4 version
%  The look of the A5 trial (space/booklet/a5-test) on an A4 page:
%    · colour cover from the unit's cover art, black-and-white inside
%    · headings in Outfit (the site's heading face), body in Lato
%    · each experiment opens with a band: number, title, ink symbol, heavy rule
%    · boxes told apart by frame and icon, never by colour:
%        I do ●  solid · We do ●● dashed · You do ●●● heavy
%        Key idea / technique ◆ thin · Above level ▲ double · Country ◎ dotted
%    · skill tags are ink tabs carrying the code (QP1, PC4 …) plus the name
%  Everything above this block still names the old colours; they are mapped to
%  ink here, so the body never needs to know which look it is printed in.
% ══════════════════════════════════════════════════════════════════════════════

\usetikzlibrary{patterns}
% ── Fonts ───────────────────────────────────────────────────────────────────
% Bold is Outfit SemiBold, the weight the site uses; Outfit has no italic,
% so a slanted one stands in.
\renewfontfamily\headingfont{Outfit}[
  Path=./fonts/, Extension=.ttf, Ligatures=TeX,
  UprightFont=Outfit-Regular, BoldFont=Outfit-SemiBold,
  ItalicFont=Outfit-Regular, ItalicFeatures={FakeSlant=0.18},
  BoldItalicFont=Outfit-SemiBold, BoldItalicFeatures={FakeSlant=0.18} ]

% ── Black and white inside ──────────────────────────────────────────────────
\definecolor{Ink}{HTML}{111111}
\colorlet{NavyBlue}{Ink}\colorlet{MidBlue}{Ink}\colorlet{TealGreen}{Ink}
\colorlet{GoldDark}{Ink}\colorlet{PurpleDark}{Ink}\colorlet{OrangeDark}{Ink}
\colorlet{TableHead}{Ink}
\colorlet{Blue}{Ink}\colorlet{Teal}{Ink}\colorlet{Amber}{Ink}\colorlet{Coral}{Ink}
\colorlet{Violet}{Ink}\colorlet{Sky}{Ink}
\colorlet{LightBlue}{white}\colorlet{LightGreen}{white}\colorlet{GoldLight}{white}
\colorlet{PurpleLight}{white}\colorlet{OrangeLight}{white}\colorlet{RowAlt}{white}
\colorlet{LightGrey}{white}\colorlet{Paper}{white}
\colorlet{DarkGrey}{black!72}\colorlet{MidGrey}{black!45}\colorlet{Line}{black!45}
\arrayrulecolor{black!60}

% ── Page style ──────────────────────────────────────────────────────────────
\fancypagestyle{fancy}{%
  \fancyhf{}%
  \renewcommand{\headrulewidth}{0.4pt}%
  \renewcommand{\headrule}{\hbox to\headwidth{\color{Ink}\leaders\hrule height \headrulewidth\hfill}}%
  \fancyhead[L]{\headingfont\footnotesize \unitname\ \textbullet\ Experiment Booklet}%
  \fancyhead[R]{\headingfont\footnotesize\leftmark}%
  \fancyfoot[C]{\headingfont\footnotesize\thepage}}
\fancypagestyle{plain}{\fancyhf{}\renewcommand{\headrulewidth}{0pt}%
  \fancyfoot[C]{\headingfont\footnotesize\thepage}}
\fancypagestyle{opener}{\fancyhf{}\renewcommand{\headrulewidth}{0pt}%
  \fancyfoot[C]{\headingfont\footnotesize\thepage}}
\pagestyle{fancy}
\markboth{}{}

\titleformat{\subsubsection}{\headingfont\normalsize\bfseries}{}{0em}{}

% ── Box system: told apart by frame and icon, never by colour ───────────────
\newcommand{\icodots}[1]{\tikz[baseline=-0.6ex]{\foreach \k in {1,...,#1}{\fill[white] (\k*1.9mm,0) circle (0.7mm);}}}
\newcommand{\icokey}{\tikz[baseline=-0.6ex]{\fill[white] (0,0.9mm)--(0.9mm,0)--(0,-0.9mm)--(-0.9mm,0)--cycle;}}
\newcommand{\icoup}{\tikz[baseline=-0.6ex]{\fill[white] (-1mm,-0.8mm)--(1mm,-0.8mm)--(0,1mm)--cycle;}}
\newcommand{\icocountry}{\tikz[baseline=-0.6ex]{\draw[white,line width=0.5pt] (0,0) circle (1mm); \fill[white] (0,0) circle (0.45mm);}}
\newcommand{\icoskill}{\tikz[baseline=-0.6ex]{\fill[white] (-0.85mm,-0.85mm) rectangle (0.85mm,0.85mm);}}
\newcommand{\tagicon}[1]{\tikz[baseline=-0.6ex]\node[fill=Ink, rounded corners=0.8mm, inner xsep=1.2mm, inner ysep=1mm]{#1};}
\tcbset{wbbase/.style={enhanced, colback=white, arc=1.2mm,
    left=3.5mm, right=3.5mm, top=4.2mm, bottom=2.6mm, before skip=11pt, after skip=8pt,
    attach boxed title to top left={xshift=3mm, yshift=-2.4mm},
    boxed title style={colback=Ink, colframe=Ink, arc=0.8mm, boxrule=0pt,
                       left=1.6mm, right=1.8mm, top=0.6mm, bottom=0.6mm},
    fonttitle=\headingfont\bfseries\footnotesize, coltitle=white}}

% I do: the worked example (solid frame, one dot)
\renewtcolorbox{workedexbox}[1]{wbbase, breakable, colframe=Ink, boxrule=1pt,
  title={\icodots{1}\ I do \textbullet\ #1}}
\renewenvironment{workedex}[1]
  {\refstepcounter{wenum}\label{we:\thewenum}\begin{workedexbox}{Worked example \thewenum\ --- #1}}
  {\end{workedexbox}}
% We do: straight after its worked example (dashed frame, two dots)
\renewtcolorbox{wedobox}[1]{wbbase, breakable=false, frame hidden,
  borderline={1.1pt}{0pt}{Ink, dash pattern=on 4pt off 2.6pt},
  title={\icodots{2}\ #1}}
\renewenvironment{wedo}[1]
  {\par\begin{wedobox}{We do \thewenum\ --- #1}\label{wd:\thewenum}\ignorespaces}
  {\end{wedobox}}
% You do: a set of practice tasks on your own (heavy frame, three dots)
\newtcolorbox{youdobox}[1][Check your understanding]{wbbase, breakable=false,
  colframe=Ink, boxrule=2.4pt, top=5mm,
  title={\icodots{3}\ You do \textbullet\ #1}}
% Key idea, technique, safety, word banks (thin frame, diamond)
\renewtcolorbox{infobox}[1]{wbbase, breakable, colframe=Ink, boxrule=0.5pt,
  title={\icokey\ #1}}
\renewtcolorbox{scaffold}[1]{wbbase, breakable, colframe=Ink, boxrule=0.5pt,
  title={\icokey\ #1}}
\renewtcolorbox{technique}[2]{wbbase, breakable, colframe=Ink, boxrule=0.5pt,
  title={\icokey\ Technique \textbullet\ #1\enspace\textbar\enspace #2 \skname{#2}}}
% Above level (double frame, triangle)
\renewtcolorbox{yr8box}{wbbase, breakable, frame hidden,
  borderline={0.6pt}{0pt}{Ink}, borderline={0.6pt}{1.6pt}{Ink},
  title={\icoup\ Going further \textbullet\ Above level}}
% Country (dotted frame, ring)
\newtcolorbox{countrybox}[1]{wbbase, breakable, frame hidden,
  borderline={1.4pt}{0pt}{Ink, line cap=round, dash pattern=on 0pt off 2.4pt},
  title={\icocountry\ #1}}
% Learning intentions and skills focus (thin frame; square for skills)
\renewtcolorbox{lipanel}{wbbase, breakable=false, colframe=Ink, boxrule=0.5pt,
  title={\icokey\ Learning intentions}}
\renewtcolorbox{skillsfocus}{wbbase, breakable=false, colframe=Ink, boxrule=1pt,
  title={\icoskill\ Skills focus \textbullet\ the science skills this experiment builds},
  before upper={\begin{multicols}{2}\raggedcolumns}, after upper={\end{multicols}}}

% ── Skill tags: an ink tab with the code, then the name ─────────────────────
\makeatletter
\renewcommand{\skc}[1]{%
  \tikz[baseline=(c.base)]\node[rounded corners=0.8mm, fill=Ink, text=white, inner xsep=1.5mm, inner ysep=0.7mm,
    font=\headingfont\bfseries\scriptsize] (c) {#1};}
\renewcommand{\sk}[1]{\mbox{\skc{#1}\,{\headingfont\bfseries\footnotesize\skname{#1}}}}
\makeatother
\renewcommand{\focus}[2]{\par\noindent\makebox[1.1cm][l]{\skc{#1}}\parbox[t]{\dimexpr\linewidth-1.1cm\relax}{\footnotesize\RaggedRight\textbf{\skname{#1}}\,---\,#2}\par\vspace{3pt}}

% ── Experiment opener band ──────────────────────────────────────────────────
% \experimentdiv{title}{summary}{lesson}: number, title and the experiment's ink
% symbol (\expsymbol{N}), then a heavy rule. Steps the section counter, so
% tasks, worked examples and figures are N.M.
\newcommand{\expword}{Experiment}\newcommand{\expcaps}{EXPERIMENT}\newcommand{\symid}{\thesection}
\renewcommand{\experimentdiv}[3]{%
  \clearpage\thispagestyle{opener}%
  \refstepcounter{section}\markboth{\expword\ \thesection\ \textbullet\ #1}{}%
  \pdfbookmark[1]{\expword\ \thesection{} — #1}{exp.\thesection}%
  \noindent\begin{minipage}[b]{\dimexpr\linewidth-46mm\relax}
    {\headingfont\bfseries\footnotesize\addfontfeatures{LetterSpace=14}\unitcaps\ \textbullet\ \expcaps}\par\vspace{1mm}
    {\headingfont\bfseries\fontsize{40}{40}\selectfont\thesection}\par\vspace{1mm}
    {\headingfont\bfseries\fontsize{19}{22}\selectfont\RaggedRight #1\par}\vspace{1mm}
    {\small\itshape #2}\par\vspace{1mm}
    {\headingfont\footnotesize Lesson #3}
  \end{minipage}\hfill
  \begin{minipage}[b]{44mm}\centering\resizebox{!}{27mm}{\expsymbol{\symid}}\end{minipage}\par
  \vspace{2mm}\noindent\rule{\linewidth}{1.6pt}\par\vspace{1mm}}
\renewcommand{\lessoncodes}[1]{{\small\itshape VC 2.0: #1}\par\vspace{2pt}}

% Part heading inside an experiment: every part starts a new page; the
% starred form follows straight on.
\RenewDocumentCommand{\exppart}{s m m}{%
  \par\IfBooleanTF{#1}{\needspace{10\baselineskip}\vspace{8pt}}{\clearpage}%
  {\headingfont\bfseries\LARGE #2}\par\vspace{1mm}\noindent\rule{\linewidth}{0.6pt}\par\vspace{2mm}}

% Key vocabulary: a heading like the trial's, then the write-in table
\renewcommand{\vocabhead}{%
  \par\vspace{4pt}{\headingfont\bfseries\Large Key vocabulary}\par\vspace{1pt}%
  \noindent{\small Copy the meaning of each word from the slide. The first one has been done for you.}\par}

% Table headings: bold Outfit on white, the grid in ink
\renewcommand{\thead}[1]{\headingfont\bfseries #1}
\renewcommand{\tickbox}{\tikz[baseline=-0.6ex]\draw[Ink, line width=0.7pt, rounded corners=1pt] (0,-0.17) rectangle (0.34,0.17);}

% Glossary rows: \glossrow{word}{label}
\newcommand{\glossrow}[2]{\rule{0pt}{5mm}\textbf{#1} & \pageref{#2} \\ \hline}

% ── Ink symbols for the experiment openers ───────────────────────────────────
% One diagrammatic idea per experiment, drawn in ink from the same 2D
% equipment shapes the booklet teaches.
\tikzset{sym/.style={line width=1.5pt, draw=Ink, line cap=round, line join=round}}
\newcommand{\expsymbol}[1]{%
  \begin{tikzpicture}[x=1cm,y=1cm]
  \ifcase#1
  \or
    % 1 Yeast: a flask with a balloon swelling on its neck
    \draw[sym] (0.6,3.0) -- (0.6,2.3) -- (0,0) -- (2.4,0) -- (1.8,2.3) -- (1.8,3.0);
    \fill[black!20] (0.32,1.2) -- (2.08,1.2) -- (2.4,0) -- (0,0) -- cycle;
    \foreach \x/\y in {0.7/0.5,1.2/0.8,1.7/0.4,1.0/0.25,1.5/0.95}{\draw[Ink, line width=1pt, fill=white] (\x,\y) circle (0.1);}
    \draw[sym] (0.6,3.0) .. controls (-0.4,3.9) and (0.2,5.0) .. (1.2,5.0) .. controls (2.2,5.0) and (2.8,3.9) .. (1.8,3.0);
  \or
    % 2 Key: a leaf and a branching key
    \draw[sym] (0,0) .. controls (0.2,1.6) and (1.0,2.4) .. (1.9,2.6) .. controls (1.8,1.4) and (1.2,0.4) .. (0,0);
    \draw[Ink, line width=1.2pt] (0,0) -- (1.6,2.2);
    \draw[sym] (3.2,2.6) -- (3.2,1.9) (2.5,1.9) -- (3.9,1.9) (2.5,1.9) -- (2.5,1.2) (3.9,1.9) -- (3.9,1.2) (2.1,1.2) -- (2.9,1.2) (2.1,1.2) -- (2.1,0.5) (2.9,1.2) -- (2.9,0.5);
    \foreach \p in {(2.1,0.4),(2.9,0.4),(3.9,1.1)}{\fill[Ink] \p circle (0.14);}
  \or
    % 3 Quadrat: a square frame with plants inside
    \draw[Ink, line width=3pt] (0,0) rectangle (3.4,3.4);
    \foreach \x/\y in {0.7/0.5,1.9/1.6,2.8/0.6,0.9/2.4,2.6/2.6}{\draw[Ink, line width=1.6pt] (\x,\y) -- (\x,\y+0.55) (\x,\y) -- (\x-0.25,\y+0.45) (\x,\y) -- (\x+0.25,\y+0.45);}
    \foreach \x/\y in {1.5/0.6,2.2/2.9}{\draw[Ink, line width=1.2pt, fill=white] (\x,\y) circle (0.22);}
  \or
    % 4 Beaks: a bird head with a long thin bill
    \draw[sym, fill=white] (1.2,1.5) circle (1.1);
    \fill[Ink] (1.55,1.85) circle (0.16);
    \draw[sym, fill=black!25] (2.2,1.9) -- (4.4,1.35) -- (2.25,1.15) -- cycle;
    \draw[sym] (0.3,0.6) .. controls (0.6,0) and (1.6,-0.1) .. (2.0,0.4);
  \or
    % 5 Populations: two out-of-step cycles on one graph
    \draw[sym, -{Stealth[length=7pt]}] (0,0) -- (4.2,0); \draw[sym, -{Stealth[length=7pt]}] (0,0) -- (0,3.2);
    \draw[Ink, line width=2.4pt] plot[smooth] coordinates {(0.1,0.8) (0.7,2.6) (1.4,1.0) (2.1,0.6) (2.8,2.5) (3.6,0.9)};
    \draw[Ink, line width=1.6pt, dashed] plot[smooth] coordinates {(0.1,0.3) (0.9,0.5) (1.5,1.6) (2.2,0.5) (3.1,0.4) (3.8,1.5)};
  \fi
  \end{tikzpicture}}
\begin{document}
\setlength{\parindent}{0pt}
\setlength{\parskip}{4pt}

% ---- Cover: the unit's cover art, full bleed, in colour ------------------
% The art is 2:3 and A4 is taller-than-wide by less, so it fills the width and
% loses a little off the bottom; the baked-in title at the top stays whole.
\pdfbookmark[1]{Cover}{cover}
\thispagestyle{empty}
\wrcover
\begin{tikzpicture}[remember picture, overlay]
  
  % Name box, top right, clear of the title.
  \node[anchor=north east, fill=white, draw=Ink, line width=0.6pt, rounded corners=1.5mm,
        inner xsep=2.6mm, inner ysep=2.2mm, text width=40mm, font=\headingfont\small]
    at ([xshift=-6mm, yshift=-6mm]current page.north east)
    {\textbf{Name}\par\vspace{5mm}\hrule\vspace{2.6mm}\textbf{Class}\par\vspace{5mm}\hrule};
  % What this book is: the cover art says Workbook, so name the booklet here.
  \node[anchor=south east, fill=white, draw=Ink, line width=0.6pt, rounded corners=1.5mm,
        inner xsep=3.2mm, inner ysep=2.4mm, font=\headingfont\bfseries\large, align=right]
    at ([xshift=-6mm, yshift=8mm]current page.south east)
    {Experiment Booklet\\[-1pt]{\normalfont\headingfont\small Experiments 1--5}};
\end{tikzpicture}
\mbox{}
\clearpage

% ---- Inside cover ---------------------------------------------------------
\thispagestyle{plain}
{\headingfont\bfseries\fontsize{26}{30}\selectfont Biodiversity and Ecosystems}\par\vspace{1mm}
{\headingfont\Large Experiment Booklet \textbullet\ Year 7}\par\vspace{1mm}
{\small\itshape VC 2.0: VC2S8U01 \textbullet\ VC2S8U04 \textbullet\ VC2S8I01--I08 \textbullet\ VC2S8H03}\par\vspace{6mm}
\renewcommand{\arraystretch}{2.2}
\begin{tabular}{@{}l p{130mm}@{}}
\textbf{Name} & \hrulefill \\ \textbf{Class} & \hrulefill \\ \textbf{Teacher} & \hrulefill \\ \textbf{Lab group} & \hrulefill \\
\end{tabular}
\renewcommand{\arraystretch}{1}

\begin{countrybox}{Acknowledgement of Country}
This booklet was created on the unceded Country of the Wurundjeri Woi-wurrung people of the Kulin Nation. We acknowledge Elders past, present and emerging. Aboriginal and Torres Strait Islander peoples have observed, tested and passed on knowledge of Country for tens of thousands of years: science in the way this booklet teaches it, by careful observation, measurement and sharing what was found.
\end{countrybox}

\begin{infobox}{How this booklet works}
This booklet holds every experiment in the unit. Bring it to every practical lesson. Each experiment opens with its key vocabulary, learning intentions and skills focus, then runs \emph{Technique} $\rightarrow$ \emph{Plan and safety} $\rightarrow$ \emph{Method} $\rightarrow$ \emph{Results} $\rightarrow$ \emph{Analysis and evaluation} $\rightarrow$ \emph{Skill check}. Each kind of box has its own frame and mark:\par\smallskip
\renewcommand{\arraystretch}{1.35}
\begin{tabular}{@{}p{32mm}l@{}}
\textbf{Technique} \tagicon{\icokey} & How to do it properly (thin frame) \\
\textbf{I do} \tagicon{\icodots{1}} & Your teacher works an example through --- you follow (solid frame) \\
\textbf{We do} \tagicon{\icodots{2}} & You try one together, with help (dashed frame) \\
\textbf{You do} \tagicon{\icodots{3}} & On your own: check your understanding (heavy frame) \\
\textbf{Going further} \tagicon{\icoup} & Above level: Year 8 extension (double frame) \\
\textbf{Skills focus} \tagicon{\icoskill} & The science skills the experiment builds \\
\end{tabular}

\smallskip
Every skill comes from the faculty \textbf{Science Skills} list (the toolkit on page \pageref{sec:toolkit}). A skill always carries the same code, for example \sk{PC4} or \sk{PA1}, so you can follow it from this unit into every other unit this year.
\end{infobox}

\clearpage
{\headingfont\bfseries\LARGE Contents}\par\vspace{1mm}\noindent\rule{\linewidth}{0.6pt}\par\vspace{2mm}
\begin{ctab}
\renewcommand{\arraystretch}{1.5}
\begin{tabular}{|L{2.8cm}|L{10.2cm}|C{1.6cm}|}
  \hline\thead{Section} & \thead{Title} & \thead{Page} \\ \hline
  Skills & \hyperlink{toolkit}{The Science Skills toolkit} & \hyperlink{toolkit}{\pageref{sec:toolkit}} \\ \hline
  Skills & \hyperlink{tracker}{My skills tracker} & \hyperlink{tracker}{\pageref{sec:tracker}} \\ \hline
  Experiment 1 & \hyperlink{exp.1}{Is Yeast Alive? (Lesson 1)} & \hyperlink{exp.1}{\pageref{sec:1}} \\ \hline
  Experiment 2 & \hyperlink{exp.2}{Key the Schoolyard (Lesson 4)} & \hyperlink{exp.2}{\pageref{sec:2}} \\ \hline
  Experiment 3 & \hyperlink{exp.3}{Quadrat Survey (Lesson 5)} & \hyperlink{exp.3}{\pageref{sec:3}} \\ \hline
  Experiment 4 & \hyperlink{exp.4}{Beaks and Food (Lesson 8)} & \hyperlink{exp.4}{\pageref{sec:4}} \\ \hline
  Experiment 5 & \hyperlink{exp.5}{The Rabbit and Fox Model (Lesson 9)} & \hyperlink{exp.5}{\pageref{sec:5}} \\ \hline
  Glossary & \hyperlink{glossary}{Key vocabulary, with the page each word first appears on} & \hyperlink{glossary}{\pageref{sec:glossary}} \\ \hline
\end{tabular}
\end{ctab}

% ══════════════════════════════════════════════════════════════════════════════
%  THE SCIENCE SKILLS TOOLKIT — faculty continuum, Years 7–8 column
% ══════════════════════════════════════════════════════════════════════════════
\clearpage
\hypertarget{toolkit}{}\pdfbookmark[1]{The Science Skills toolkit}{toolkitbm}
\markboth{The Science Skills toolkit}{}{\headingfont\bfseries\LARGE The Science Skills toolkit}\label{sec:toolkit}\par\vspace{1mm}\noindent\rule{\linewidth}{0.6pt}\par
These are the skills scientists use in every investigation. They come from the Science Inquiry strand of the Victorian Curriculum, grouped the way our faculty teaches them. Each experiment in this booklet tells you which skills it builds.

\newcommand{\tkhead}[2]{\multicolumn{3}{|l|}{\cellcolor{black!10}\rule{0pt}{4.6mm}\headingfont\bfseries #2} \\ \hline}
\newcommand{\tkrow}[1]{\skc{#1}\enspace{\small\bfseries\skname{#1}} & \skdesc{#1} & {\small\skcode{#1}} \\ \hline}
\renewcommand{\arraystretch}{1.35}
\begin{ctab}
\begin{longtable}{|W{5.0cm}|W{9.0cm}|Z{1.7cm}|}
  \hline\thead{Skill} & \thead{What it means} & \thead{Code} \\ \hline
  \endhead
  \tkhead{Teal}{Questioning and predicting}
  \tkrow{QP1}
  \tkrow{QP2}
  \tkhead{Blue}{Planning and conducting}
  \tkrow{PC1}
  \tkrow{PC2}
  \tkrow{PC3}
  \tkrow{PC4}
  \tkrow{PC5}
  \tkrow{PC6}
  \tkrow{PC7}
  \tkrow{PC8}
  \tkhead{GoldDark}{Processing, modelling and analysing}
  \tkrow{PA1}
  \tkrow{PA2}
  \tkrow{PA3}
  \tkrow{PA4}
  \tkrow{PA5}
  \tkrow{PA6}
  \tkrow{PA7}
  \tkrow{PA8}
  \tkhead{OrangeDark}{Evaluating}
  \tkrow{EV1}
  \tkrow{EV2}
  \tkrow{EV3}
  \tkrow{EV4}
  \tkrow{EV5}
  \tkhead{PurpleDark}{Communicating}
  \tkrow{CO1}
  \tkrow{CO2}
\end{longtable}
\end{ctab}
\renewcommand{\arraystretch}{1}

\begin{infobox}{Writing like a scientist \textbar\ CO1 \skname{CO1}}
Scientists write in the \textbf{third person}: say what was done, not who did it. Write \emph{``The water was heated for 10 minutes''}, not \emph{``I heated the water for 10 minutes''}. Keep sentences short, spell scientific words correctly, and always give a unit with a measurement.
\end{infobox}

% ══════════════════════════════════════════════════════════════════════════════
%  SKILLS TRACKER
% ══════════════════════════════════════════════════════════════════════════════
\clearpage
\hypertarget{tracker}{}\pdfbookmark[1]{My skills tracker}{trackerbm}
\markboth{My skills tracker}{}{\headingfont\bfseries\LARGE My skills tracker}\label{sec:tracker}\par\vspace{1mm}\noindent\rule{\linewidth}{0.6pt}\par
A dot \textbullet{} shows the experiments that build each skill. After each experiment, write how confident you are with each of its skills in its column: \textbf{1}~not yet, \textbf{2}~nearly, \textbf{3}~got it.

\newcommand{\dotc}{{\Large\textbullet}}
\newcommand{\trow}[6]{\rule[-0.36cm]{0pt}{0.92cm}\skc{#1}\enspace{\small\skname{#1}} & #2 & #3 & #4 & #5 & #6 & \\ \hline}
\begin{ctab}
\begin{tabular}{|W{5.0cm}|*{5}{Z{1.3cm}|}Z{1.8cm}|}
  \hline\thead{Skill} & \thead{Exp 1} & \thead{Exp 2} & \thead{Exp 3} & \thead{Exp 4} & \thead{Exp 5} & \thead{End of unit} \\ \hline
  \trow{QP1}{\dotc}{}{\dotc}{\dotc}{}
  \trow{PC1}{}{}{\dotc}{}{}
  \trow{PC2}{\dotc}{}{}{\dotc}{}
  \trow{PC5}{}{\dotc}{}{}{}
  \trow{PC6}{\dotc}{\dotc}{\dotc}{}{}
  \trow{PA1}{\dotc}{}{\dotc}{\dotc}{\dotc}
  \trow{PA2}{}{}{\dotc}{\dotc}{\dotc}
  \trow{PA3}{}{}{\dotc}{}{}
  \trow{PA4}{}{\dotc}{}{}{\dotc}
  \trow{PA5}{}{}{}{}{\dotc}
  \trow{PA6}{}{}{}{\dotc}{}
  \trow{PA8}{}{}{}{\dotc}{\dotc}
  \trow{EV1}{\dotc}{}{}{}{}
  \trow{EV3}{}{\dotc}{\dotc}{}{}
  \trow{CO2}{}{\dotc}{}{}{}
\end{tabular}
\end{ctab}
{\small\color{DarkGrey} The other skills are built in the experiment booklets for other units.}

% ══════════════════════════════════════════════════════════════════════════════
%  EXPERIMENT 1 — IS YEAST ALIVE?                         (Lesson 1, T2 Week 9)
% ══════════════════════════════════════════════════════════════════════════════
\experimentdiv{Is Yeast Alive?}{Testing for a life process \textbullet{} A control \textbullet{} Measuring a balloon over time}{1}\label{sec:1}
\lessoncodes{VC2S8U01 \textbullet{} VC2S8I01, I02, I04, I06}

\vocabhead
\begin{ctab}
\begin{tabular}{|SW{4.2cm}|W{11.6cm}|}
  \hline\hd & \thead{Word} & \thead{Definition (copy from the slide)} \\ \hline
  & \scaf{\textbf{Life process}\label{voc:life-process}} & \scaf{Something every living thing does: movement, respiration, sensitivity, growth, reproduction, excretion, nutrition (MRS GREN).} \\ \hline
  & \textbf{Respiration}\label{voc:respiration} & \\ \hline
  & \textbf{Micro-organism}\label{voc:micro-organism} & \\ \hline
  & \textbf{Control}\label{voc:control} & \\ \hline
  & \textbf{Circumference}\label{voc:circumference} & \\ \hline
\end{tabular}
\end{ctab}

\begin{lipanel}
\inquiryli{Write a hypothesis, use a control and record measurements over time.}
\understandli{Use evidence of a life process to decide whether yeast is living.}
\sctitle
\begin{itemize}[leftmargin=1.4em, itemsep=0pt, topsep=1pt]
  \item I can write a hypothesis that names both variables.
  \item I can explain what my control flask shows.
  \item I can use my results to argue that yeast is alive.
\end{itemize}
\end{lipanel}

\begin{skillsfocus}
\focus{QP1}{Write a testable hypothesis.}
\focus{PC2}{Name the variables and the control.}
\focus{PC6}{Measure each balloon every five minutes.}
\focus{PA1}{Record measurements in a THUD table.}
\focus{EV1}{Say whether the results support the hypothesis.}
\end{skillsfocus}

% ─── TECHNIQUE ────────────────────────────────────────────────────────────────
\exppart{Technique spotlight}{Blue}

\begin{technique}{What a control is for}{PC2}
A \textbf{control} is a set-up that is the same as the others except for the one thing you are testing. If the yeast with sugar blows up its balloon and the yeast \emph{without} sugar does not, the sugar made the difference, not the warm water or the flask. Without a control you cannot tell what caused the result.
\end{technique}

\begin{technique}{Measuring a round object}{PC6}
\begin{minipage}[c]{0.62\linewidth}
Wrap a piece of string once around the widest part of the balloon. Pinch the string where it meets its end, then lay it straight along a ruler. Read the length to the nearest 0.5\,cm. This is the \textbf{circumference}. Measure the same way, at the same place, every time.
\end{minipage}\hfill
\begin{minipage}[c]{0.34\linewidth}\centering
\begin{tikzpicture}[scale=0.7]
  \draw[eq] (0,0) .. controls (-1.3,0.6) and (-1.2,2.4) .. (0,2.6) .. controls (1.2,2.4) and (1.3,0.6) .. (0,0);
  \draw[eq] (-0.15,0) -- (0.15,0) -- (0.1,-0.25) -- (-0.1,-0.25) -- cycle;
  \draw[Ink, line width=1.2pt, dashed] (-1.02,1.4) .. controls (-0.5,1.15) and (0.5,1.15) .. (1.02,1.4);
  \eqlab{(1.0,1.4)}{(2.0,2.0)}{string}
\end{tikzpicture}
\end{minipage}
\end{technique}

\begin{workedex}{Writing the hypothesis}
\emph{If sugar is added to the yeast, the circumference of the balloon will increase.}

\smallskip
With a reason: \emph{\dots because the yeast uses the sugar for respiration and gives off carbon dioxide gas.}
\end{workedex}

\begin{wedo}{Writing the hypothesis}
Write a hypothesis for flask C, which has sugar and yeast but \emph{cold} water.
\answerbox{3cm}
\end{wedo}

% ─── PLAN AND SAFETY / METHOD ────────────────────────────────────────────────
\exppart{Plan, safety and method}{Teal}

\textbf{Aim:} To investigate the effect of changing the food and temperature of yeast on the circumference of a balloon over its flask.

\task Complete the variables table. \skr{PC2}
\begin{ctab}
\begin{tabular}{|SW{4.4cm}|W{10.2cm}|}
  \hline\hd & \thead{Variable} & \thead{In this experiment} \\ \hline
  & \scaf{Independent} & \scaf{What is in the flask: sugar or no sugar; warm or cold water} \\ \hline
  & Dependent & \\ \hline
  & Controlled & \\ \hline
  & Controlled & \\ \hline
  & Control flask & \\ \hline
\end{tabular}
\end{ctab}

\needspace{8\baselineskip}
\task Write your hypothesis for flask A. \skr{QP1}
\answerbox{2.5cm}

\subsubsection*{Equipment}
3 conical flasks (250\,mL) labelled A, B and C \textbullet{} 3 balloons \textbullet{} 3 teaspoons of dried yeast \textbullet{} 2 teaspoons of sugar \textbullet{} warm water (about 40\,\textdegree C) and cold water \textbullet{} 100\,mL measuring cylinder \textbullet{} funnel \textbullet{} string and a ruler \textbullet{} stopwatch

\begin{center}
\begin{tikzpicture}[scale=0.72]
  \foreach \x/\l/\c in {0/A/{warm water, sugar, yeast},4.2/B/{warm water, yeast},8.4/C/{cold water, sugar, yeast}}{
    \begin{scope}[shift={(\x,0)}]
      \eqFlask
      \fill[black!12] (-0.12,0.03) -- (0.1,0.7) -- (1.7,0.7) -- (1.72,0.03) -- cycle;
      \draw[eq] (0.55,2.95) .. controls (0.1,3.6) and (0.3,4.5) .. (0.8,4.6) .. controls (1.3,4.5) and (1.5,3.6) .. (1.05,2.95);
      \node[font=\headingfont\bfseries] at (0.8,1.5) {\l};
      \node[eqtext, align=center, text width=3.2cm] at (0.8,-0.55) {\c};
    \end{scope}}
\end{tikzpicture}
\expfigcap{Flask B has no sugar; flask C is cold.}
\end{center}

\needspace{12\baselineskip}
\task Complete the risk assessment. \skr{PC3}
\begin{ctab}
\begin{tabular}{|SW{4.0cm}|W{5.2cm}|W{5.4cm}|}
  \hline\hd & \thead{Risk (what could go wrong)} & \thead{Prevention} & \thead{Responding (if it happens)} \\ \hline
  & \scaf{A glass flask breaks} & \scaf{Keep flasks in the middle of the bench; carry one at a time} & \scaf{Do not pick it up; tell the teacher, who sweeps it up} \\ \hline
  & Warm water spills & & \\ \hline
  & Someone is allergic to latex balloons & & \\ \hline
\end{tabular}
\end{ctab}

\subsubsection*{Method}
\begin{enumerate}[itemsep=2pt]
  \item Using the funnel, put 1 teaspoon of yeast into each flask, and 1 teaspoon of sugar into flasks A and C.
  \item Add 100\,mL of warm water to flasks A and B, and 100\,mL of cold water to flask C. Swirl each gently.
  \item Stretch a balloon over the neck of each flask straight away, and start the stopwatch.
  \item Measure the circumference of each balloon at 0, 5, 10, 15 and 20 minutes. Record each in Table 1.1.
  \item Pour the contents down the sink with plenty of water, and rinse the flasks.
\end{enumerate}

% ─── RESULTS ─────────────────────────────────────────────────────────────────
\exppart{Results}{Amber}

\subsubsection*{Table 1.1: Circumference of each balloon over time \skr{PA1}}
\begin{ctab}
\begin{tabular}{|SW{2.6cm}|C{2.2cm}|C{2.2cm}|C{2.2cm}|C{2.2cm}|C{2.2cm}|}
  \hline\hd & \thead{Flask} & \thead{0 min (cm)} & \thead{5 min (cm)} & \thead{10 min (cm)} & \thead{15 min (cm)} & \thead{20 min (cm)} \\ \hline
  & \scaf{Example} & \scaf{6.0} & \scaf{8.5} & \scaf{12.0} & \scaf{15.5} & \scaf{18.0} \\ \hline
  & A & & & & & \\ \hline
  & B & & & & & \\ \hline
  & C & & & & & \\ \hline
\end{tabular}
\end{ctab}

\needspace{12\baselineskip}
\task Describe what you saw in each flask, apart from the balloon. \skr{PC6}
\answerbox{4cm}

% ─── ANALYSIS ────────────────────────────────────────────────────────────────
\exppart{Analysis and evaluation}{Coral}

\task Describe the trend in flask A. Use numbers from Table 1.1. \skr{PA5}
\answerbox{3.5cm}

\needspace{10\baselineskip}
\task Flask B is a control. What does it tell you about the result in flask A?
\answerbox{4cm}

\needspace{10\baselineskip}
\task Explain the result in flask C.
\answerbox{3.5cm}

\needspace{12\baselineskip}
\task Which life processes (MRS GREN) did your experiment give evidence for? For each, say what the evidence was.
\answerbox{5cm}

\needspace{12\baselineskip}
\task Write a conclusion: state the trend, back it with numbers from your table, and say whether your hypothesis was supported. \skr{EV1}

\scaf{The yeast with sugar in warm water made the balloon grow from \dots\,cm to \dots\,cm. This supports / does not support the hypothesis.}
\answerbox{5cm}

\begin{yr8box}
\begin{keepext}
\exttask A student boils yeast for five minutes, cools it, and repeats flask A with it. Predict what will happen to the balloon, and explain what this would show about dead and living things.
\answerbox{5cm}
\end{keepext}
\end{yr8box}

% ─── SKILL CHECK ─────────────────────────────────────────────────────────────
\exppart{Skill check}{NavyBlue}

\begin{skillcheck}
\scrow{QP1}{I can write a hypothesis that names both variables.}
\scrow{PC2}{I can explain what a control is for.}
\scrow{PC6}{I measured each balloon the same way every five minutes.}
\scrow{PA1}{My table has a title, headings, units in the headings and matching decimals.}
\scrow{EV1}{I can say whether my results supported the hypothesis.}
\end{skillcheck}

\task Which skill will you focus on in Experiment 2? Explain what you will do differently.
\answerbox{5cm}
% ══════════════════════════════════════════════════════════════════════════════
%  EXPERIMENT 2 — KEY THE SCHOOLYARD                     (Lesson 4, T3 Week 1)
% ══════════════════════════════════════════════════════════════════════════════
\experimentdiv{Key the Schoolyard}{Observable features \textbullet{} Building a dichotomous key \textbullet{} Testing it on someone else}{4}\label{sec:2}
\lessoncodes{VC2S8U01 \textbullet{} VC2S8I02, I04, I06, I08}

\vocabhead
\begin{ctab}
\begin{tabular}{|SW{4.2cm}|W{11.6cm}|}
  \hline\hd & \thead{Word} & \thead{Definition (copy from the slide)} \\ \hline
  & \scaf{\textbf{Dichotomous key}\label{voc:dichotomous-key}} & \scaf{A tool that identifies an organism through a series of questions, each with two possible answers.} \\ \hline
  & \textbf{Feature}\label{voc:feature} & \\ \hline
  & \textbf{Observable}\label{voc:observable} & \\ \hline
  & \textbf{Specimen}\label{voc:specimen} & \\ \hline
  & \textbf{Margin}\label{voc:margin} & \\ \hline
\end{tabular}
\end{ctab}

\begin{lipanel}
\inquiryli{Build a model (a key) from careful observations, and test whether someone else can use it.}
\understandli{Classify organisms by their observable features.}
\sctitle
\begin{itemize}[leftmargin=1.4em, itemsep=0pt, topsep=1pt]
  \item I can write yes/no questions about features anyone can see.
  \item I can draw my key as a branching diagram.
  \item I can find the questions in my key that confused someone, and fix them.
\end{itemize}
\end{lipanel}

\begin{skillsfocus}
\focus{PC5}{Collect specimens without harming living things.}
\focus{PC6}{Observe each specimen closely.}
\focus{PA4}{Build a key: a model for sorting.}
\focus{CO2}{Draw the key as a branching diagram.}
\focus{EV3}{Find and fix the questions that caused errors.}
\end{skillsfocus}

% ─── TECHNIQUE ────────────────────────────────────────────────────────────────
\exppart{Technique spotlight}{Blue}

\begin{technique}{Writing a good question for a key}{PA4}
Each question splits the specimens into two groups, using a feature anyone can \textbf{see or measure}:
\begin{itemize}[leftmargin=1.4em, itemsep=1pt, topsep=3pt]
  \item \textbf{Good:} \emph{Is the leaf longer than 5\,cm?} \quad \emph{Does the leaf margin have teeth?} \quad \emph{Is the leaf split into smaller leaflets?}
  \item \textbf{Poor:} \emph{Is the leaf pretty?} (opinion) \quad \emph{Is it big?} (big compared with what?) \quad \emph{Is it a gum leaf?} (that is the answer, not a feature)
\end{itemize}
Split the specimens roughly in half each time and the key stays short.
\end{technique}

\begin{workedex}{A key for four leaves}
\begin{center}
\begin{tikzpicture}[font=\small, node distance=6mm,
  q/.style={draw=Ink, line width=0.8pt, rounded corners=1.5pt, align=center, inner sep=3pt, text width=3.6cm},
  a/.style={draw=Ink, line width=1.4pt, align=center, inner sep=3pt, font=\headingfont\bfseries\small}]
  \node[q] (q1) {Is the leaf split into leaflets?};
  \node[q, below left=8mm and -6mm of q1] (q2) {Does the margin have teeth?};
  \node[a, below right=8mm and 4mm of q1] (w) {Wattle};
  \node[a, below left=8mm and -4mm of q2] (r) {Rose};
  \node[a, below right=8mm and -4mm of q2] (g) {Gum};
  \draw[-{Stealth[length=4mm]}, Ink] (q1) -- node[left, font=\scriptsize] {no} (q2);
  \draw[-{Stealth[length=4mm]}, Ink] (q1) -- node[right, font=\scriptsize] {yes} (w);
  \draw[-{Stealth[length=4mm]}, Ink] (q2) -- node[left, font=\scriptsize] {yes} (r);
  \draw[-{Stealth[length=4mm]}, Ink] (q2) -- node[right, font=\scriptsize] {no} (g);
\end{tikzpicture}
\end{center}
\end{workedex}

\begin{wedo}{Fix the question}
A student's key asks \emph{Is the leaf a nice shape?} Explain what is wrong with it and write a better question.
\answerbox{3.5cm}
\end{wedo}

% ─── PLAN AND SAFETY / METHOD ────────────────────────────────────────────────
\exppart{Plan, safety and method}{Teal}

\textbf{Aim:} To build a dichotomous key that identifies eight schoolyard leaves, and test whether another group can use it.

\begin{infobox}{Collecting responsibly}
Collect \textbf{fallen} leaves wherever you can, and take only one leaf from a living plant. Leave animals where you find them: if you want a minibeast in your key, draw or photograph it and put it back. Scientists collect only what they need.
\end{infobox}

\needspace{10\baselineskip}
\task Complete the risk assessment. \skr{PC3}
\begin{ctab}
\begin{tabular}{|SW{4.0cm}|W{5.2cm}|W{5.4cm}|}
  \hline\hd & \thead{Risk (what could go wrong)} & \thead{Prevention} & \thead{Responding (if it happens)} \\ \hline
  & \scaf{A plant irritates the skin} & \scaf{Do not touch plants with milky sap or hairs; wash hands afterwards} & \scaf{Wash the skin with water; tell the teacher} \\ \hline
  & A bee or ant sting & & \\ \hline
  & Sunburn & & \\ \hline
\end{tabular}
\end{ctab}

\subsubsection*{Equipment}
8 envelopes or a tray \textbullet{} hand lens \textbullet{} ruler \textbullet{} pencil \textbullet{} sticky labels numbered 1 to 8

\subsubsection*{Method}
\begin{enumerate}[itemsep=2pt]
  \item Collect eight different leaves from the schoolyard. Label them 1 to 8 and write where each came from.
  \item Record the features of each leaf in Table 2.1.
  \item Choose a first question that splits the leaves into two groups. Keep asking questions until every leaf is on its own.
  \item Draw your key as a branching diagram. Give each leaf a name: its real name if you know it, or a made-up name if you do not.
  \item Swap leaves and keys with another group. Use their key to name each of their leaves, then check with them.
\end{enumerate}

% ─── RESULTS ─────────────────────────────────────────────────────────────────
\exppart{Results}{Amber}

\subsubsection*{Table 2.1: Features of eight schoolyard leaves \skr{PC6}}
\begin{ctab}
\begin{tabular}{|SW{1.4cm}|W{2.6cm}|C{2.2cm}|W{2.6cm}|C{2.4cm}|W{2.6cm}|}
  \hline\hd & \thead{Leaf} & \thead{Where found} & \thead{Length (cm)} & \thead{Margin (smooth / toothed)} & \thead{Leaflets (yes / no)} & \thead{Other feature} \\ \hline
  & \scaf{e.g.} & \scaf{Oval fence} & \scaf{11.5} & \scaf{Smooth} & \scaf{No} & \scaf{Curved, waxy} \\ \hline
  & 1 & & & & & \\ \hline
  & 2 & & & & & \\ \hline
  & 3 & & & & & \\ \hline
  & 4 & & & & & \\ \hline
  & 5 & & & & & \\ \hline
  & 6 & & & & & \\ \hline
  & 7 & & & & & \\ \hline
  & 8 & & & & & \\ \hline
\end{tabular}
\end{ctab}

\needspace{20\baselineskip}
\task Draw your dichotomous key as a branching diagram. \skr{CO2}
\drawbox{12cm}{}

\needspace{12\baselineskip}
\subsubsection*{Table 2.2: Testing another group's key \skr{EV3}}
\begin{ctab}
\begin{tabular}{|SW{1.4cm}|W{5.2cm}|W{5.2cm}|C{2.0cm}|}
  \hline\hd & \thead{Leaf} & \thead{Name from their key} & \thead{Correct name (from them)} & \thead{Correct?} \\ \hline
  & 1 & & & \\ \hline
  & 2 & & & \\ \hline
  & 3 & & & \\ \hline
  & 4 & & & \\ \hline
  & 5 & & & \\ \hline
  & 6 & & & \\ \hline
  & 7 & & & \\ \hline
  & 8 & & & \\ \hline
\end{tabular}
\end{ctab}

% ─── ANALYSIS ────────────────────────────────────────────────────────────────
\exppart{Analysis and evaluation}{Coral}

\task How many leaves did the other group identify correctly with your key? Which question caused a mistake? \skr{EV3}
\answerbox{4cm}

\needspace{10\baselineskip}
\task Rewrite that question so it could not be misread.
\answerbox{3cm}

\needspace{10\baselineskip}
\task A key is a \emph{model} for sorting living things. Give one thing your key does well, and one thing it could not do (for example, a leaf from a ninth plant). \skr{PA8}
\answerbox{4cm}

\needspace{10\baselineskip}
\task Why should a key use features you can see, rather than where the leaf was found or what season it is?
\answerbox{4cm}

\begin{yr8box}
\begin{keepext}
\exttask Field guides to Australian birds often use a key based on size, bill shape and colour. Explain why colour can be a poor feature for a key, using an example such as a young bird, a male and a female.
\answerbox{5cm}
\end{keepext}
\end{yr8box}

% ─── SKILL CHECK ─────────────────────────────────────────────────────────────
\exppart{Skill check}{NavyBlue}

\begin{skillcheck}
\scrow{PC5}{I collected specimens without harming living things.}
\scrow{PC6}{I recorded the features of each leaf carefully.}
\scrow{PA4}{I built a key that separates every specimen.}
\scrow{CO2}{I drew my key as a clear branching diagram.}
\scrow{EV3}{I found the question that caused errors and fixed it.}
\end{skillcheck}

\task Which skill will you focus on in Experiment 3? Explain what you will do differently.
\answerbox{5cm}
% ══════════════════════════════════════════════════════════════════════════════
%  EXPERIMENT 3 — QUADRAT SURVEY                         (Lesson 5, T3 Week 1)
% ══════════════════════════════════════════════════════════════════════════════
\experimentdiv{Quadrat Survey}{Sampling with quadrats \textbullet{} Averages from repeat samples \textbullet{} Comparing two sites}{5}\label{sec:3}
\lessoncodes{VC2S8U04 \textbullet{} VC2S8I01, I02, I04, I06}

\vocabhead
\begin{ctab}
\begin{tabular}{|SW{4.2cm}|W{11.6cm}|}
  \hline\hd & \thead{Word} & \thead{Definition (copy from the slide)} \\ \hline
  & \scaf{\textbf{Quadrat}\label{voc:quadrat}} & \scaf{A square frame, often 1\,m by 1\,m, used to sample the organisms in one small area.} \\ \hline
  & \textbf{Biodiversity}\label{voc:biodiversity} & \\ \hline
  & \textbf{Species richness}\label{voc:species-richness} & \\ \hline
  & \textbf{Sample}\label{voc:sample} & \\ \hline
  & \textbf{Random sampling}\label{voc:random-sampling} & \\ \hline
\end{tabular}
\end{ctab}

\begin{lipanel}
\inquiryli{Collect data with quadrats, average repeat samples and compare two sites in a column graph.}
\understandli{Compare the biodiversity of two habitats and explain the difference using abiotic and biotic factors.}
\sctitle
\begin{itemize}[leftmargin=1.4em, itemsep=0pt, topsep=1pt]
  \item I can place quadrats randomly and count what is inside.
  \item I can calculate an average from three quadrats.
  \item I can explain why one site has more species than the other.
\end{itemize}
\end{lipanel}

\begin{skillsfocus}
\focus{QP1}{Write a hypothesis about the two sites.}
\focus{PC1}{Follow the sampling method exactly.}
\focus{PC6}{Count every species in each quadrat.}
\focus{PA1}{Record counts in a THUD table.}
\focus{PA2}{Draw a column graph of the averages.}
\focus{PA3}{Calculate averages.}
\focus{EV3}{Explain how sampling can give errors.}
\end{skillsfocus}

% ─── TECHNIQUE ────────────────────────────────────────────────────────────────
\exppart{Technique spotlight}{Blue}

\begin{technique}{Why we sample}{PC1}
You cannot count every plant and animal on the whole oval, so you count a few small areas (\textbf{samples}) and use them to stand for the whole site. Two rules keep a sample fair:
\begin{itemize}[leftmargin=1.4em, itemsep=1pt, topsep=3pt]
  \item \textbf{Random:} you do not choose the interesting spots. Stand still, close your eyes and toss the quadrat gently over your shoulder.
  \item \textbf{Repeated:} one quadrat could land on an unusual patch. Take at least three and average them.
\end{itemize}
\end{technique}

\begin{technique}{Counting species richness}{PC6}
\begin{minipage}[c]{0.6\linewidth}
\textbf{Species richness} is the number of \emph{different} kinds of living things, not the total number of living things. Ten clover plants and one ant is a richness of 2. You do not need names: ``broad grass'', ``thin grass'', ``small black ant'' are fine, as long as you use the same name at both sites.
\end{minipage}\hfill
\begin{minipage}[c]{0.36\linewidth}\centering
\begin{tikzpicture}[scale=0.55]
  \draw[Ink, line width=1.6pt] (0,0) rectangle (4,4);
  \foreach \x/\y in {0.6/0.7,1.4/2.9,2.6/1.2,3.3/3.2,0.8/3.4,2.0/2.2}{\draw[Ink, line width=0.8pt] (\x,\y) -- (\x,\y+0.45) (\x,\y) -- (\x-0.2,\y+0.35) (\x,\y) -- (\x+0.2,\y+0.35);}
  \foreach \x/\y in {1.2/1.0,3.2/0.6,2.8/2.7}{\draw[Ink, fill=white, line width=0.8pt] (\x,\y) circle (0.22); \foreach \a in {0,72,...,288}{\draw[Ink] (\x,\y) ++(\a:0.22) -- ++(\a:0.14);}}
  \fill[Ink] (3.6,1.8) ellipse (0.12 and 0.07);
  \node[font=\scriptsize, anchor=north] at (2,-0.1) {1\,m};
\end{tikzpicture}
\end{minipage}
\end{technique}

\begin{workedex}{Richness and the average}
Three quadrats on a garden bed had species richness 7, 9 and 8.

\smallskip
Average $= (7 + 9 + 8) \div 3 = 24 \div 3 = 8$ species.
\end{workedex}

\begin{wedo}{Counting richness}
The quadrat in the picture above holds grass, a flowering weed and one ant. What is its species richness?
\answerbox{2cm}
\end{wedo}

% ─── PLAN AND SAFETY / METHOD ────────────────────────────────────────────────
\exppart{Plan, safety and method}{Teal}

\textbf{Aim:} To investigate the effect of changing the site (the mown oval or a garden bed) on species richness.

\task Write a hypothesis. \skr{QP1}
\answerbox{2.5cm}

\needspace{12\baselineskip}
\task Complete the variables table. \skr{PC2}
\begin{ctab}
\begin{tabular}{|SW{4.4cm}|W{10.2cm}|}
  \hline\hd & \thead{Variable} & \thead{In this experiment} \\ \hline
  & \scaf{Independent} & \scaf{The site: mown oval or garden bed} \\ \hline
  & Dependent & \\ \hline
  & Controlled & \\ \hline
  & Controlled & \\ \hline
\end{tabular}
\end{ctab}

\needspace{10\baselineskip}
\task Complete the risk assessment. \skr{PC3}
\begin{ctab}
\begin{tabular}{|SW{4.0cm}|W{5.2cm}|W{5.4cm}|}
  \hline\hd & \thead{Risk (what could go wrong)} & \thead{Prevention} & \thead{Responding (if it happens)} \\ \hline
  & \scaf{Stung by a bee or bull ant} & \scaf{Look before you kneel; do not put hands under logs} & \scaf{Tell the teacher; anyone allergic uses their EpiPen plan} \\ \hline
  & A thrown quadrat hits someone & & \\ \hline
  & Sunburn & & \\ \hline
\end{tabular}
\end{ctab}

\subsubsection*{Equipment}
1\,m $\times$ 1\,m quadrat (or four metre rulers) \textbullet{} hand lens \textbullet{} clipboard and pencil \textbullet{} hat and sunscreen

\subsubsection*{Method}
\begin{enumerate}[itemsep=2pt]
  \item At the first site, stand still with your back to the area, check no one is behind you, and toss the quadrat gently over your shoulder.
  \item Search the quadrat for 3 minutes, including under leaves. Record each different kind of plant and animal, and count how many kinds there are.
  \item Repeat steps 1 and 2 twice more at the same site.
  \item Repeat steps 1 to 3 at the second site.
  \item Calculate the average species richness for each site.
\end{enumerate}

% ─── RESULTS ─────────────────────────────────────────────────────────────────
\exppart{Results}{Amber}

\subsubsection*{Table 3.1: Species richness at two sites \skr{PA1}}
\begin{ctab}
\begin{tabular}{|SW{3.4cm}|C{2.6cm}|C{2.6cm}|C{2.6cm}|C{3.0cm}|}
  \hline\hd & \thead{Site} & \thead{Quadrat 1 (species)} & \thead{Quadrat 2 (species)} & \thead{Quadrat 3 (species)} & \thead{Average (species)} \\ \hline
  & \scaf{Example: path edge} & \scaf{4} & \scaf{6} & \scaf{5} & \scaf{5} \\ \hline
  & Mown oval & & & & \\ \hline
  & Garden bed & & & & \\ \hline
\end{tabular}
\end{ctab}

\task Show your working for one average. \skr{PA3}
\workbox{2.5cm}

\needspace{12\baselineskip}
\task List the kinds of living things you found at each site.
\begin{ctab}
\begin{tabular}{|W{7.4cm}|W{7.4cm}|}
  \hline\thead{Mown oval} & \thead{Garden bed} \\ \hline
  \rule{0pt}{4cm} & \\ \hline
\end{tabular}
\end{ctab}

\needspace{18\baselineskip}
\task Draw a column graph of the average species richness at each site. Give it a title and label both axes with units. \skr{PA2}
\graphgrid{8cm}

% ─── ANALYSIS ────────────────────────────────────────────────────────────────
\exppart{Analysis and evaluation}{Coral}

\task Which site had greater biodiversity? Use numbers from Table 3.1. \skr{PA5}
\answerbox{3.5cm}

\needspace{12\baselineskip}
\task Explain the difference between the sites. Use at least one abiotic factor (such as shade, water or mowing) and one biotic factor.
\answerbox{5cm}

\needspace{12\baselineskip}
\task Give two reasons your averages might not match the true richness of each site, and how a better survey could fix each. \skr{EV3}
\begin{ctab}
\begin{tabular}{|SW{7.2cm}|W{7.8cm}|}
  \hline\hd & \thead{Source of error} & \thead{How to reduce it} \\ \hline
  & \scaf{Only three quadrats per site} & \scaf{Take ten quadrats per site and average them} \\ \hline
  & & \\ \hline
  & & \\ \hline
\end{tabular}
\end{ctab}

\needspace{10\baselineskip}
\task Write a conclusion: the trend, the numbers that back it, and whether your hypothesis was supported. \skr{EV1}
\answerbox{4.5cm}

\begin{yr8box}
\begin{keepext}
\exttask Species richness counts kinds, not how many of each. Garden A has 9 kinds, each with about 10 individuals. Garden B has 9 kinds, but 90\% of the individuals are one grass. Which garden is more diverse? Explain why richness alone does not show this.
\answerbox{5cm}
\end{keepext}
\end{yr8box}

% ─── SKILL CHECK ─────────────────────────────────────────────────────────────
\exppart{Skill check}{NavyBlue}

\begin{skillcheck}
\scrow{QP1}{I can write a hypothesis that names both variables.}
\scrow{PC1}{I followed the sampling method exactly, including random placement.}
\scrow{PC6}{I counted every kind of living thing in each quadrat.}
\scrow{PA1}{My table has a title, headings, units and matching decimals.}
\scrow{PA2}{I drew a column graph with a title and labelled axes.}
\scrow{PA3}{I calculated averages correctly.}
\scrow{EV3}{I can explain how sampling can give errors.}
\end{skillcheck}

\task Which skill will you focus on in Experiment 4? Explain what you will do differently.
\answerbox{5cm}
% ══════════════════════════════════════════════════════════════════════════════
%  EXPERIMENT 4 — BEAKS AND FOOD                         (Lesson 8, T3 Week 3)
% ══════════════════════════════════════════════════════════════════════════════
\experimentdiv{Beaks and Food}{Modelling an adaptation \textbullet{} Fair trials \textbullet{} Correlation and cause}{8}\label{sec:4}
\lessoncodes{VC2S8U04 \textbullet{} VC2S8I01, I02, I04, I05}

\vocabhead
\begin{ctab}
\begin{tabular}{|SW{4.2cm}|W{11.6cm}|}
  \hline\hd & \thead{Word} & \thead{Definition (copy from the slide)} \\ \hline
  & \scaf{\textbf{Adaptation}\label{voc:adaptation}} & \scaf{A feature or behaviour that helps an organism survive and reproduce in its habitat.} \\ \hline
  & \textbf{Structural adaptation}\label{voc:structural-adaptation} & \\ \hline
  & \textbf{Niche}\label{voc:niche} & \\ \hline
  & \textbf{Correlation}\label{voc:correlation} & \\ \hline
  & \textbf{Cause}\label{voc:cause} & \\ \hline
\end{tabular}
\end{ctab}

\begin{lipanel}
\inquiryli{Run fair repeated trials, graph the results and decide whether a pattern shows cause.}
\understandli{Explain how the shape of a beak is a structural adaptation to a bird's food.}
\sctitle
\begin{itemize}[leftmargin=1.4em, itemsep=0pt, topsep=1pt]
  \item I can keep the trials fair: same time, same food, same rules.
  \item I can draw a column graph of the averages.
  \item I can link each ``beak'' to a real Australian bird and its food.
\end{itemize}
\end{lipanel}

\begin{skillsfocus}
\focus{QP1}{Write a hypothesis for one food.}
\focus{PC2}{Keep every trial fair.}
\focus{PA1}{Record counts in a THUD table.}
\focus{PA2}{Draw a column graph.}
\focus{PA6}{Decide whether the pattern shows cause.}
\focus{PA8}{Say what the model leaves out.}
\end{skillsfocus}

% ─── TECHNIQUE ────────────────────────────────────────────────────────────────
\exppart{Technique spotlight}{Blue}

\begin{technique}{Correlation or cause?}{PA6}
Two things are \textbf{correlated} when they change together: as one goes up, the other goes up (or down). That alone does not show that one \textbf{causes} the other. You can claim cause only when you changed one thing yourself and kept everything else the same, as in a fair test.

\smallskip
\emph{Ice-cream sales and shark sightings rise together.} Correlated, but ice-cream does not cause sharks: hot weather causes both.
\end{technique}

\begin{workedex}{Correlation or cause?}
\textbf{Claim:} ``The spoon picked up more rice than the tweezers, so a spoon-shaped beak causes better rice-eating.''

\smallskip
In this experiment the beak was the only thing changed (same rice, same 30 seconds, same student), so the difference was caused by the beak shape \emph{in the model}. Whether it applies to real birds needs evidence from real birds.
\end{workedex}

\begin{wedo}{Correlation or cause?}
Students who eat breakfast get higher test marks. Does breakfast cause higher marks? Give one other thing that could explain the pattern.
\answerbox{3.5cm}
\end{wedo}

% ─── PLAN AND SAFETY / METHOD ────────────────────────────────────────────────
\exppart{Plan, safety and method}{Teal}

\textbf{Aim:} To investigate the effect of changing the shape of a ``beak'' on the amount of each food it can collect in 30 seconds.

\begin{center}
\begin{tikzpicture}[scale=0.62, font=\scriptsize]
  % tweezers
  \begin{scope}[shift={(0,0)}]
    \draw[eq] (0,0) -- (1.6,2.6) (0.35,0) -- (1.7,2.5); \draw[eq] (1.6,2.6) -- (1.7,2.5);
    \node at (0.85,-0.45) {tweezers};
  \end{scope}
  % peg
  \begin{scope}[shift={(3.0,0)}]
    \draw[eq] (0,0) rectangle (0.32,2.6); \draw[eq] (0.42,0) rectangle (0.74,2.6); \draw[eq] (-0.1,1.4) rectangle (0.84,1.8);
    \node at (0.37,-0.45) {peg};
  \end{scope}
  % spoon
  \begin{scope}[shift={(5.6,0)}]
    \draw[eq] (0.5,0.55) ellipse (0.5 and 0.55); \draw[eq] (0.5,1.1) -- (0.5,2.6);
    \node at (0.5,-0.45) {spoon};
  \end{scope}
  % straw
  \begin{scope}[shift={(8.0,0)}]
    \draw[eq] (0,0) rectangle (0.3,2.6);
    \node at (0.15,-0.45) {straw};
  \end{scope}
  % foods
  \begin{scope}[shift={(10.4,0)}]
    \draw[eq] (0,0.2) .. controls (0,-0.1) and (1.6,-0.1) .. (1.6,0.2); \draw[eq] (0,0.2) -- (1.6,0.2);
    \foreach \x in {0.3,0.6,0.9,1.2}{\fill[Ink] (\x,0.35) ellipse (0.09 and 0.04);}
    \node at (0.8,-0.45) {rice};
    \draw[eq] (2.2,0.2) .. controls (2.2,-0.1) and (3.8,-0.1) .. (3.8,0.2); \draw[eq] (2.2,0.2) -- (3.8,0.2);
    \foreach \x in {2.5,2.9,3.4}{\draw[Ink, fill=white] (\x,0.42) circle (0.17);}
    \node at (3.0,-0.45) {marbles};
    \draw[eq] (4.4,0) rectangle (5.6,1.0); \fill[black!15] (4.42,0.02) rectangle (5.58,0.7);
    \node at (5.0,-0.45) {``nectar''};
  \end{scope}
\end{tikzpicture}
\expfigcap{Four ``beaks'' and three foods.}
\end{center}

\task Choose one food. Write a hypothesis about which beak will collect the most of it. \skr{QP1}
\answerbox{2.5cm}

\needspace{12\baselineskip}
\task Complete the variables table. \skr{PC2}
\begin{ctab}
\begin{tabular}{|SW{4.4cm}|W{10.2cm}|}
  \hline\hd & \thead{Variable} & \thead{In this experiment} \\ \hline
  & \scaf{Independent} & \scaf{The beak (tweezers, peg, spoon or straw)} \\ \hline
  & Dependent & \\ \hline
  & Controlled & \\ \hline
  & Controlled & \\ \hline
\end{tabular}
\end{ctab}

\subsubsection*{Equipment}
Tweezers, clothes peg, plastic spoon, straw \textbullet{} rice grains, marbles, a cup of coloured water (``nectar'') \textbullet{} 3 shallow trays \textbullet{} a cup as the ``stomach'' \textbullet{} 10\,mL measuring cylinder \textbullet{} stopwatch

\begin{infobox}{Safety}
Marbles on the floor are a slip hazard: pick up any that fall straight away. The straw is used like a dropper, never sucked: nothing goes in your mouth in the lab.
\end{infobox}

\subsubsection*{Method}
\begin{enumerate}[itemsep=2pt]
  \item Spread 50 rice grains in one tray and 20 marbles in another. Put 50\,mL of ``nectar'' in a cup.
  \item The ``bird'' holds the beak in one hand and the stomach cup in the other. The timer says go.
  \item For 30 seconds, move as much food as possible into the stomach cup, using only the beak. For nectar, dip the straw in, cover its top with a finger, lift it out and let the nectar run into the stomach cup.
  \item Count the rice grains or marbles, or measure the nectar in mL. Reset the tray.
  \item Do three trials for each beak with each food, swapping the ``bird'' each time.
\end{enumerate}

% ─── RESULTS ─────────────────────────────────────────────────────────────────
\exppart{Results}{Amber}

\subsubsection*{Table 4.1: Average food collected in 30 seconds (three trials) \skr{PA1}}
\begin{ctab}
\begin{tabular}{|SW{3.2cm}|C{3.4cm}|C{3.4cm}|C{3.8cm}|}
  \hline\hd & \thead{Beak} & \thead{Rice (grains)} & \thead{Marbles (marbles)} & \thead{Nectar (mL)} \\ \hline
  & \scaf{Example: chopsticks} & \scaf{11.3} & \scaf{4.0} & \scaf{0.0} \\ \hline
  & Tweezers & & & \\ \hline
  & Peg & & & \\ \hline
  & Spoon & & & \\ \hline
  & Straw & & & \\ \hline
\end{tabular}
\end{ctab}
{\small\color{DarkGrey}Record each trial in your workbook or on the back page, then write the averages here to one decimal place.}

\needspace{18\baselineskip}
\task Draw a column graph of the average rice collected by each beak. Give it a title and label both axes with units. \skr{PA2}
\graphgrid{8cm}

% ─── ANALYSIS ────────────────────────────────────────────────────────────────
\exppart{Analysis and evaluation}{Coral}

\task Which beak was best for each food? Use numbers from Table 4.1. \skr{PA5}
\answerbox{4cm}

\needspace{12\baselineskip}
\task Match each beak to a real bird and its food. \skr{PA8}
\begin{ctab}
\begin{tabular}{|SW{3.0cm}|W{5.6cm}|W{6.0cm}|}
  \hline\hd & \thead{Beak} & \thead{Real bird} & \thead{Its food} \\ \hline
  & \scaf{Straw} & \scaf{New Holland honeyeater} & \scaf{Nectar from flowers such as grevillea} \\ \hline
  & Tweezers & & \\ \hline
  & Peg & & \\ \hline
  & Spoon & & \\ \hline
\end{tabular}
\end{ctab}

\needspace{12\baselineskip}
\task Did the beak shape \emph{cause} the differences in your model? Explain using how the trials were run. \skr{PA6}
\answerbox{4.5cm}

\needspace{10\baselineskip}
\task Give two ways this model is different from a real bird feeding. \skr{PA8}
\answerbox{4cm}

\needspace{10\baselineskip}
\task Write a conclusion for the food you chose: the trend, numbers that back it, and whether your hypothesis was supported. \skr{EV1}
\answerbox{4.5cm}

\begin{yr8box}
\begin{keepext}
\exttask On the Galápagos Islands, a drought killed most plants with small soft seeds. The next year, finches with deeper beaks were more common. Explain this using the idea of adaptation, and say why the finches did not ``grow'' deeper beaks on purpose.
\answerbox{5cm}
\end{keepext}
\end{yr8box}

% ─── SKILL CHECK ─────────────────────────────────────────────────────────────
\exppart{Skill check}{NavyBlue}

\begin{skillcheck}
\scrow{QP1}{I can write a hypothesis that names both variables.}
\scrow{PC2}{I kept every trial fair.}
\scrow{PA1}{My table has a title, headings, units and matching decimals.}
\scrow{PA2}{I drew a column graph with a title and labelled axes.}
\scrow{PA6}{I can tell correlation from cause.}
\scrow{PA8}{I can say what a model shows well and what it leaves out.}
\end{skillcheck}

\task Which skill will you focus on in Experiment 5? Explain what you will do differently.
\answerbox{5cm}
% ══════════════════════════════════════════════════════════════════════════════
%  EXPERIMENT 5 — THE RABBIT AND FOX MODEL               (Lesson 9, T3 Week 3)
% ══════════════════════════════════════════════════════════════════════════════
\experimentdiv{The Rabbit and Fox Model}{Simulating two populations \textbullet{} A line graph over generations \textbullet{} Judging a model}{9}\label{sec:5}
\lessoncodes{VC2S8U04, VC2S8H03 \textbullet{} VC2S8I04, I05, I06}

\vocabhead
\begin{ctab}
\begin{tabular}{|SW{4.2cm}|W{11.6cm}|}
  \hline\hd & \thead{Word} & \thead{Definition (copy from the slide)} \\ \hline
  & \scaf{\textbf{Population}\label{voc:population}} & \scaf{All the organisms of one species living in an area at the same time.} \\ \hline
  & \textbf{Predator}\label{voc:predator} & \\ \hline
  & \textbf{Prey}\label{voc:prey} & \\ \hline
  & \textbf{Introduced species}\label{voc:introduced-species} & \\ \hline
  & \textbf{Simulation}\label{voc:simulation} & \\ \hline
\end{tabular}
\end{ctab}

\begin{lipanel}
\inquiryli{Run a simulation, graph two populations over time and evaluate the model.}
\understandli{Explain how predator and prey populations depend on each other, using two introduced species.}
\sctitle
\begin{itemize}[leftmargin=1.4em, itemsep=0pt, topsep=1pt]
  \item I can follow the rules of the simulation exactly for 15 generations.
  \item I can draw both populations on one line graph and describe the pattern.
  \item I can say what the model shows well and what it leaves out.
\end{itemize}
\end{lipanel}

\begin{skillsfocus}
\focus{PA1}{Record both populations each generation.}
\focus{PA2}{Draw two lines on one graph, with a key.}
\focus{PA4}{Use a physical model of two populations.}
\focus{PA5}{Describe the pattern in the two lines.}
\focus{PA8}{Evaluate the model against the real world.}
\end{skillsfocus}

% ─── TECHNIQUE ────────────────────────────────────────────────────────────────
\exppart{Technique spotlight}{Blue}

\begin{technique}{Two lines on one graph}{PA2}
When two things are measured at the same times, draw them on the same axes so they can be compared:
\begin{itemize}[leftmargin=1.4em, itemsep=1pt, topsep=3pt]
  \item Time (generation) goes on the horizontal axis.
  \item Use one scale on the vertical axis that fits the larger population.
  \item Use a different mark for each: $\times$ joined by a solid line for rabbits, $\circ$ joined by a dashed line for foxes.
  \item Add a \textbf{key} that says which line is which.
\end{itemize}
Here each point is an exact count, not a measurement with error, so joining the points is the right choice.
\end{technique}

\begin{technique}{Judging a model}{PA8}
A model is simpler than the real thing on purpose. To judge it, ask:
\begin{enumerate}[leftmargin=1.6em, itemsep=1pt, topsep=3pt]
  \item What does it show \textbf{well}?
  \item What does it \textbf{leave out}?
  \item How could it be made \textbf{better}?
\end{enumerate}
\end{technique}

\begin{wedo}{Judging a model}
A globe is a model of Earth. Give one thing it shows well and one thing it leaves out.
\answerbox{3cm}
\end{wedo}

% ─── PLAN AND METHOD ─────────────────────────────────────────────────────────
\exppart{The simulation}{Teal}

\begin{readingbox}{Background --- rabbits and foxes in Australia}
European rabbits were released in Victoria in 1859 and spread across most of the continent within 50 years. Red foxes were released in the 1870s, partly to control them. Both are \textbf{introduced species}: foxes now prey on native animals such as bandicoots as well as rabbits. In this simulation, paper squares stand for rabbits and a card stands for each fox.
\end{readingbox}

\textbf{Aim:} To model how the numbers of rabbits (prey) and foxes (predators) change over 15 generations.

\subsubsection*{Equipment}
An A3 sheet as the ``field'' \textbullet{} 200 paper squares (2\,cm) as rabbits \textbullet{} 20 card squares (6\,cm) as foxes

\subsubsection*{The rules}
\begin{enumerate}[itemsep=2pt]
  \item \textbf{Start:} spread 5 rabbits on the field. Have 1 fox ready.
  \item \textbf{Hunt:} drop each fox onto the field from 30\,cm above it. Every rabbit the fox card touches is eaten: take it off.
  \item \textbf{Foxes:} a fox that caught 3 or more rabbits survives and has one cub (add one fox for next time). A fox that caught fewer than 3 dies.
  \item \textbf{Rabbits:} every rabbit left on the field has one baby: add one rabbit for each. Spread them out. (The field holds at most 200.)
  \item \textbf{Record} the number of rabbits and foxes in Table 5.1. That is one generation.
  \item If all the foxes die, 1 new fox moves in next generation. If all the rabbits are eaten, 3 new rabbits move in.
\end{enumerate}

% ─── RESULTS ─────────────────────────────────────────────────────────────────
\exppart{Results}{Amber}

\subsubsection*{Table 5.1: Rabbits and foxes over 15 generations \skr{PA1}}
\begin{ctab}
\begin{tabular}{|SW{2.6cm}|C{2.2cm}|C{2.2cm}||C{2.6cm}|C{2.2cm}|C{2.2cm}|}
  \hline\hd & \thead{Generation} & \thead{Rabbits} & \thead{Foxes} & \thead{Generation} & \thead{Rabbits} & \thead{Foxes} \\ \hline
  & \scaf{0} & \scaf{5} & \scaf{1} & 8 & & \\ \hline
  & 1 & & & 9 & & \\ \hline
  & 2 & & & 10 & & \\ \hline
  & 3 & & & 11 & & \\ \hline
  & 4 & & & 12 & & \\ \hline
  & 5 & & & 13 & & \\ \hline
  & 6 & & & 14 & & \\ \hline
  & 7 & & & 15 & & \\ \hline
\end{tabular}
\end{ctab}

\needspace{20\baselineskip}
\task Draw both populations on one line graph. Join the points, use a different line for each animal, and add a key and a title. \skr{PA2}
\graphgrid{10cm}

% ─── ANALYSIS ────────────────────────────────────────────────────────────────
\exppart{Analysis and evaluation}{Coral}

\task Describe the pattern of the rabbit line. Use numbers from your table. \skr{PA5}
\answerbox{4cm}

\needspace{10\baselineskip}
\task What happened to the foxes soon after the rabbits increased? Explain why.
\answerbox{4cm}

\needspace{10\baselineskip}
\task What happened to the rabbits after the foxes increased? Explain why.
\answerbox{4cm}

\needspace{14\baselineskip}
\task Judge the model. \skr{PA8}
\begin{ctab}
\begin{tabular}{|W{4.4cm}|W{10.4cm}|}
  \hline\thead{Question} & \thead{Your answer} \\ \hline
  What it shows well & \rule{0pt}{1.6cm} \\ \hline
  What it leaves out (give two) & \rule{0pt}{2.2cm} \\ \hline
  How to make it better & \rule{0pt}{1.6cm} \\ \hline
\end{tabular}
\end{ctab}

\needspace{12\baselineskip}
\task Fox baiting is used in parts of Victoria to protect native animals. Use your model to predict what might happen to the rabbit population if all the foxes in an area were removed, and suggest one problem this could cause for farmers.
\answerbox{5cm}

\begin{yr8box}
\begin{keepext}
\exttask Calicivirus was released in Australia in 1996 and killed many rabbits. Add a rule to the simulation that models a disease, run five more generations with it, and describe what changes for both populations.
\answerbox{6cm}
\end{keepext}
\end{yr8box}

% ─── SKILL CHECK ─────────────────────────────────────────────────────────────
\exppart{Skill check}{NavyBlue}

\begin{skillcheck}
\scrow{PA1}{I recorded both populations every generation.}
\scrow{PA2}{I drew two lines on one graph with a key and a title.}
\scrow{PA4}{I followed the rules of the model exactly.}
\scrow{PA5}{I can describe the pattern in both lines, quoting numbers.}
\scrow{PA8}{I can say what a model shows well, what it leaves out and how to improve it.}
\end{skillcheck}

\task Look back over all five experiments. Which skill have you improved most this unit? Give an example.
\answerbox{5cm}

\clearpage
\hypertarget{glossary}{}\pdfbookmark[1]{Glossary}{glossarybm}
\markboth{Glossary}{}{\headingfont\bfseries\LARGE Glossary}\label{sec:glossary}\par\vspace{1mm}\noindent\rule{\linewidth}{0.6pt}\par
Every key vocabulary word in this booklet, and the page where you wrote what it means.\par
\begin{ctab}\renewcommand{\arraystretch}{1.5}
\begin{tabular}{|L{6cm}|C{2cm}|}
\hline\thead{Word} & \thead{Page} \\ \hline
\glossrow{Adaptation}{voc:adaptation}
\glossrow{Biodiversity}{voc:biodiversity}
\glossrow{Cause}{voc:cause}
\glossrow{Circumference}{voc:circumference}
\glossrow{Control}{voc:control}
\glossrow{Correlation}{voc:correlation}
\glossrow{Dichotomous key}{voc:dichotomous-key}
\glossrow{Feature}{voc:feature}
\glossrow{Introduced species}{voc:introduced-species}
\glossrow{Life process}{voc:life-process}
\glossrow{Margin}{voc:margin}
\glossrow{Micro-organism}{voc:micro-organism}
\glossrow{Niche}{voc:niche}
\glossrow{Observable}{voc:observable}
\glossrow{Population}{voc:population}
\glossrow{Predator}{voc:predator}
\glossrow{Prey}{voc:prey}
\glossrow{Quadrat}{voc:quadrat}
\glossrow{Random sampling}{voc:random-sampling}
\glossrow{Respiration}{voc:respiration}
\glossrow{Sample}{voc:sample}
\glossrow{Simulation}{voc:simulation}
\glossrow{Species richness}{voc:species-richness}
\glossrow{Specimen}{voc:specimen}
\glossrow{Structural adaptation}{voc:structural-adaptation}
\end{tabular}
\end{ctab}

\end{document}
